\documentclass[10pt,a4paper]{article}
\usepackage[margin=2.1cm]{geometry}
\usepackage{fontspec}
\usepackage{xeCJK}
\setCJKmainfont{Noto Serif CJK KR}
\setCJKsansfont{Noto Sans CJK KR}
\setCJKmonofont{Noto Sans Mono CJK KR}
\xeCJKsetup{CJKspace=true}
\usepackage{amsmath,amssymb,mathtools}
\usepackage{booktabs,tabularx,multirow,array,longtable}
\usepackage{enumitem}
\usepackage{xcolor}
\usepackage{caption}
\usepackage{tikz}
\usetikzlibrary{positioning,arrows.meta,fit,backgrounds,calc,shapes.geometric,decorations.pathreplacing}
\usepackage[colorlinks=true,linkcolor=blue!50!black,citecolor=green!40!black,urlcolor=blue!60!black]{hyperref}
\linespread{1.15}
\setlength{\parskip}{3pt}
\setlist{itemsep=1pt,topsep=2pt}

% ---------- notation macros ----------
\newcommand{\G}{\mathcal{G}}
\newcommand{\KS}{\mathcal{K}_S}
\newcommand{\Ont}{\mathcal{O}}
\newcommand{\Ctx}{\mathcal{C}}
\newcommand{\Tcmd}{\mathcal{T}_{\mathrm{cmd}}}
\newcommand{\Treq}{\mathcal{T}_{\mathrm{req}}}
\newcommand{\Tst}{\mathcal{T}_{\mathrm{state}}}
\newcommand{\Tadv}{\mathcal{T}_{\mathrm{adv}}}
\newcommand{\Yset}{\mathcal{Y}}
\newcommand{\Vdom}{\mathcal{V}_{\mathrm{dom}}}
\newcommand{\KR}{\mathcal{K}_R}
\newcommand{\Cfg}{\Theta_{\mathrm{kg}}}
\newcommand{\Ivl}{\mathbb{I}}
\newcommand{\ATCO}{\textsf{ATCO}}
\newcommand{\PILOT}{\textsf{PILOT}}
\newcommand{\cmd}[1]{\textsf{\small #1}}
\newcommand{\ch}[1]{\textbf{\textcolor{red!60!black}{#1}}}
\newcommand{\md}[1]{\textbf{\textcolor{blue!60!black}{#1}}}

\definecolor{perc}{RGB}{225,238,250}
\definecolor{comp}{RGB}{232,245,228}
\definecolor{proj}{RGB}{252,240,222}
\definecolor{res}{RGB}{240,236,248}

\title{\textbf{미국 ATC 무선교신 기반 사건 중심 Temporal Knowledge Graph}\\[2pt]
\large 개발 시스템 문제정의 및 모듈 설계 (System Module Specification)}
\author{Team 6}
\date{2026-10-14 \quad v2: 구현 반영 (\S\ref{sec:impl}) \quad (문제정의.md 및 문제정의\_및\_KG구축방안\_v2.md의 후속 문서)}

\begin{document}
\maketitle

\begin{abstract}
\noindent 이 문서는 \emph{문제정의.md}에서 세운 연구 질문 RQ1--RQ4와, 공유 샘플(248초)에서 \emph{측정·관찰}한 기술적 문제 \ch{C1}--\ch{C15}를 개발 시스템의 입출력 명세로 옮긴다.
시스템은 구축 모듈 \md{M0}--\md{M6}과 RQ4 진단 모듈 \md{M7}로 구성된다. \md{M1} STT는 송신 분할(\md{M1a})과 전사(\md{M1b})로, \md{M2} 지시 추출은 화자 역할 판정(\md{M2a})와 의미 추출(\md{M2b})로 이루어진다. STT를 MOSS-ATC(30분 통녹음의 전사·송신 경계·화자 번호를 한 번에 출력)로 쓰면 \md{M1a}의 음성 구간 분할은 모델이 대신한다. 구축 데이터는 HF LiveATC 원본 아카이브(KATL)로 확정했다(\S\ref{sec:data}).
각 모듈마다 (i) \textbf{해결하는 기술적 문제}, (ii) \textbf{입력·출력의 형식 정의}, (iii) \textbf{방법과 적합성 근거}(선행연구)를 명시한다.
마지막에 문제 $\leftrightarrow$ 모듈 \textbf{추적성 행렬}(\S\ref{sec:trace})과 샘플 실행 예(\S\ref{sec:example})를 제시하여, 문제정의와 제안 방법이 빠짐없이 1:1 이상 대응됨을 보인다.
\end{abstract}

\tableofcontents
\newpage

%==========================================================
\section*{구현 반영 사항 (v2)}\label{sec:impl}
\addcontentsline{toc}{section}{구현 반영 사항 (v2)}
%==========================================================
아래는 설계(본문 \S1--\S11)를 실제로 구현하면서 바뀐 점과, KATL Ground 5.84시간에 대한 실행 결과다. 본문의 수식 명세는 그대로 유효하고, 바뀐 부분은 이 절이 우선한다. 코드는 \texttt{src/} 아래에 있다.

\begin{center}\small
\begin{tabularx}{\textwidth}{@{}l>{\raggedright\arraybackslash}p{4.0cm}>{\raggedright\arraybackslash}X@{}}
\toprule
모듈 & 설계 & 구현 (이유)\\\midrule
M0 & FAA NASR 또는 OurAirports & 활주로·주파수는 OurAirports, 유도로 102개는 OpenStreetMap에서 가져온다 (NASR에는 유도로 이름이 없다). 결과 \texttt{data/static\_kb/katl.json}은 \textbf{이름 목록}이며 좌표·연결 정보는 없다. \texttt{kb/build\_static\_kb.py}\\
M1a & 에너지 문턱 + VAD & faster-whisper에 내장된 Silero VAD만 쓴다(최소 무음 300\,ms). 원본 아카이브는 디지털 무음이 0\%이기 때문이다. 30분 파일끼리 약 2분 겹치는 부분은 UTC 기준으로 중복 제거한다.\\
M1b & 맥락 프롬프트 & 같다. 직전 파일의 활성 콜사인을 다음 파일의 프롬프트에 넣는다(KG$\to$STT 되먹임). 도시명은 \texttt{municipality}에서 가져온다(``Atlanta Ground'').\\
M2a & 콜사인 위치 규칙 & \textbf{화자 단위 LLM 판정}으로 바꿨다. MOSS가 준 화자 번호별로 송신을 모아 LLM이 화자마다 역할 하나를 정한다(\texttt{speaker\_roles.json}). LLM 답이 없으면 그 화자의 송신에 규칙 판정을 다수결한다. Whisper 경로(화자 번호 없음)는 송신마다 규칙으로 판정하고 역할이 바뀌는 곳에서 자른다. 규칙만으로는 역할 미정이 51\%였다.\\
M2b & JSON Schema 제약 디코딩 & \textbf{사후 검증 + 하이브리드}로 바꿨다. (1) LLM이 JSON을 출력한다(Qwen3-4B-Instruct, 로컬). (2) 온톨로지 유형, 대상 형태, 원문 근거, \textbf{명령별 어휘 단서}(\texttt{ontology.json: lexical\_cues})로 검증한다. (3) LLM이 놓친 유형은 규칙 문법으로 보충한다. 어휘 단서가 없으면 지상 주파수에서 HEADING·CLEARED\_FOR\_TAKEOFF 같은 환각이 125·33건 나왔고, 단서 적용 후 3·0건으로 줄었다. 원 출력은 저장해 두어 재검증(\texttt{--revalidate})할 수 있다.\\
M3 & 콜사인 해소, KB 연결 & 같다. 추가로 \cmd{ABBREVIATED}(번호만 부름)를 최근 활성 항공편과 대조해 해소한다. \texttt{Alpha+번호} 오인식은 같은 번호의 활성 항공편으로 교정한다(\cmd{ASR\_CORRECTED}). ``Delta''는 음성 알파벳과 겹치므로 예외로 처리한다.\\
M4 & Neo4j 저장, SHACL 검증 & \textbf{파이썬 이중 시간 그래프(JSON)}가 주 저장소이고, Neo4j는 CSV + \texttt{import.cypher} \emph{내보내기}로 지원한다. 제약은 SHACL 대신 코드로 검사한다(데모·평가를 설치 없이 실행하기 위해). 개념이 겹치지 않는 조종사 잡담은 복창으로 연결하지 않는다.\\
M5 & $\rho_1$--$\rho_7$ & 같다. 경로 비교는 \emph{일부 생략은 허용, 치환은 불일치}로 정의했다. $\rho_4$에서 동시 활주로 횡단(CROSS--CROSS)은 정상 운영이라 제외한다.\\
M6 & 시점 질의 & 같다. FastAPI 웹 데모(\texttt{src/demo/}): strip board, 시간 스크러버, 시점 스냅샷, 2열 교신 기록, 사건 그래프, 근거 원음 재생.\\
M7 & 진단 하네스 & E1--E4 스크립트(\texttt{src/evaluation/}). E2 LLM 판정은 추출과 다른 모델(Qwen3.5-9B)을 쓴다.\\
온톨로지 & SESAR 온톨로지 재사용 & \textbf{DLR 공식 사양 v3.00(설명서·명령 표 PDF)} \cite{dlr2025}과 \textbf{FAA JO 7110.65BB} \cite{faa7110}를 입력으로 \texttt{kb/build\_ontology.py}가 생성한다. 사람이 정한 대응(20개 종류 $\leftrightarrow$ DLR 표의 행, FAA 조항·표준 문구)을 코드가 PDF 원문과 대조하고, 하나라도 원문에 없으면 생성을 멈춘다. 기계 판독 공식 파일(OWL/JSON)은 없다. 18개 종류가 DLR 행에 대응하고, \cmd{MONITOR}(FAA §2-1-17)와 \cmd{REQUEST}(DLR의 Reason 속성)는 대응 행이 없다. 미국 적용 근거는 NASA--DLR 실증(명령 종류 손실 2.6\%, 값 포함 4.6\%) \cite{helmke2025loss}.\\
지식 파일 & 온톨로지 $\Ont$ 하나 & 성격별로 세 파일로 나눴다. \texttt{configs/ontology.json}(어휘·구조, DLR 근거), \texttt{configs/rules.json}(규칙 베이스 $\KR$: 복창 규칙·지시 대체·활주로 점유, FAA 근거와 지시 의미), \texttt{configs/kg\_build.json}(구축 설정 $\Cfg$: 복창 시간창 20\,s, 최대 유효시간). 항공사 호출부호는 \texttt{configs/airlines.json}. RQ4에서는 온톨로지를 그대로 두고 $\KR$의 복창 규칙만 지역 규정으로 바꾼다.\\
\bottomrule
\end{tabularx}
\end{center}

\paragraph{실행 결과 (\texttt{results/katl\_gnd}).}
\begin{center}\small
\begin{tabular}{@{}lrr@{}}
\toprule
지표 & LLM 하이브리드 & 규칙 기준선\\\midrule
송신 턴 (M2a) & 1,775 & 1,775\\
역할 미정 (UNKNOWN) & 3 & 873\\
항공편 노드 & 513 & 431\\
관제 지시 & 1,088 & 441\\
복창 링크 (MATCH/PARTIAL/MISMATCH) & 105 (71/14/20) & 69 (41/15/13)\\
신뢰도 HIGH/MEDIUM/LOW & 87/697/304 & 51/261/129\\
경보 & 312 & 240\\
E1 역량 질문 (답변·근거) & 8/8 & 8/8 (CQ5 빈 답)\\
E3 대조군 오경보 & 0\% (n=71) & 0\% (n=41)\\
E3 콜사인 숫자 변경 / 누락 & 95.5\% / 91.2\% & 97.6\% / 97.2\%\\
E3 유도로 치환 & 56.1\% & 90.0\%\\
처리 시간 (ASR / M2b) & 7.5분 / 17분 & 7.5분 / 1초 미만\\
\bottomrule
\end{tabular}
\end{center}
E2(근거 정밀도)는 LLM 판정 기준 60\%(잠정, n=60)다. 판정 모델의 오류가 확인되어 사람 라벨이 필요하다.

\newpage

%==========================================================
\section{시스템 수준 문제정의}
%==========================================================
\subsection{입력}
\begin{itemize}
\item \textbf{교신 음성}: $a \in \mathbb{R}^{N}$, 단채널, $f_s = 16\,\mathrm{kHz}$ (원본 mp3를 재표본화). 녹음 시작 시각을 $\tau_0$라 두고, 표본 $n$의 시각은 $t = \tau_0 + n/f_s$.
\item \textbf{녹음 맥락}: $\Ctx = (\alpha_{\mathrm{apt}}, \phi_{\mathrm{pos}}, \tau_0)$. 공항 ICAO 코드(예: \cmd{KSFO}), 관제 위치·주파수 $\phi_{\mathrm{pos}} \in \{\cmd{TWR},\cmd{GND},\cmd{APP},\dots\}$, 시작 시각. 샘플처럼 여러 공항 음원이 섞인 경우에는 클립 단위로 $\Ctx$를 부여한다 (\ch{C9}).
\item \textbf{사전 지식(정적 자원)}:
  \begin{itemize}
  \item 온톨로지 $\Ont = (\mathcal{T}, \mathcal{R}, \Sigma)$: 개념 유형 $\mathcal{T}$, 관계 유형 $\mathcal{R}$, 무결성 제약 $\Sigma$.
  \item 정적 도메인 KB $\KS = (V_S, E_S)$: 공항·활주로·유도로·관제시설·항공사.
  \item 규칙 베이스 $\KR = (\Treq^{\rho},\ \mathrm{sup},\ \mathcal{T}_{\mathrm{rwy}})$: 지역별 복창 필수 유형, 지시 대체 관계, 활주로 점유 유형.
  \item 구축 설정 $\Cfg = (\Delta_{\mathrm{rb}},\ T_{\max})$: 복창 허용 시간창, 유형별 최대 유효시간. 도메인 지식이 아니라 녹음 누락에 대비한 데이터 처리 설정이다.
  \end{itemize}
\end{itemize}

\subsection{출력}
\begin{itemize}
\item \textbf{사건 중심 Temporal KG}
\begin{equation}
\G = (V,\; E,\; \lambda,\; \tau,\; \mathrm{prov},\; \mathrm{conf}),
\end{equation}
  \begin{itemize}
  \item $\lambda: V \cup E \to \mathcal{T}\cup\mathcal{R}$: 유형 함수
  \item $\tau: V\cup E \to \Ivl$: 유효 구간. $\Ivl = \{[t_1,t_2) : t_1 \le t_2 \le \infty\}$
  \item $\mathrm{prov}: V\cup E \to 2^{S}$: 출처 송신 집합
  \item $\mathrm{conf}: V\cup E \to [0,1]$: 신뢰도
  \end{itemize}
\item \textbf{이상 상황 집합}
\begin{equation}
\Yset = \{(y_k,\; \pi_k,\; \gamma_k)\}_{k}, \quad y_k \in \mathcal{L}_{\mathrm{anom}},\;\; \pi_k \subseteq \G \text{ (근거 부분그래프)},\;\; \gamma_k\in[0,1],
\end{equation}
  \begin{itemize}
  \item 라벨 집합:
  \begin{align*}\mathcal{L}_{\mathrm{anom}} = \{&\cmd{MISMATCH},\cmd{PARTIAL},\cmd{NO\_READBACK},\cmd{NO\_CALLSIGN},\\ &\cmd{RWY\_OVERLAP},\cmd{POST\_HANDOFF},\cmd{KB\_VIOLATION}\}\end{align*}
  \item 정상 판정 \cmd{MATCH}는 $\G$의 \cmd{READBACK\_OF} 간선 속성으로만 남는다.
  \end{itemize}
\item \textbf{상황 질의 함수}: $Q: \G \times \mathbb{R} \times \mathcal{Q} \to \mathrm{Ans}$. 시각 $t$와 질의 템플릿 $q\in\mathcal{Q}$를 받는다.
\end{itemize}

\subsection{시스템 함수와 만족해야 할 조건}
\begin{equation}
F:\; (a,\, \Ctx;\; \Ont,\, \KS,\, \KR,\, \Cfg) \;\longmapsto\; (\G,\, \Yset,\, Q)
\end{equation}
$F$는 다음을 만족해야 한다.
\begin{enumerate}[label=(P\arabic*)]
\item \textbf{스키마 적합성}: $\G \models \Sigma$. 모든 노드·간선이 $\Ont$의 유형과 제약을 따른다 (RQ1).
\item \textbf{출처 보존}: $\forall x \in V\cup E \setminus V_S:\; \mathrm{prov}(x) \neq \emptyset$. 모든 동적 사실은 근거 송신을 가진다 (설명 가능성, \ch{C12}).
\item \textbf{시간 정합성}: $\forall x:\; \tau(x)=[t_1,t_2) \Rightarrow t_1 \ge \min_{s\in\mathrm{prov}(x)} t^{s}$. 사실은 근거 발화보다 먼저 성립할 수 없다 (\ch{C10}).
\item \textbf{정적 KB 연결}: 대상 개체 언급은 $V_S$의 노드로 연결되거나, 연결되지 않으면 명시적으로 $\bot$로 표시한다 (RQ2).
\item \textbf{탐지 가능성}: $\Yset = \mathcal{R}(\G)$. 규칙 집합 $\mathcal{R}$을 $\G$에 적용한 결과로만 이상을 도출한다 (RQ3).
\end{enumerate}

\subsection{표기 요약}
\begin{center}\small
\begin{tabular}{@{}ll@{\qquad}ll@{}}
\toprule
기호 & 의미 & 기호 & 의미\\\midrule
$s_i=[t^s_i,t^e_i)$ & $i$번째 송신(transmission) 구간 & $x_i$ & 송신 $i$의 음성 클립\\
$w_i$ & 전사 단어열 & $c_i$ & ASR 신뢰도 벡터\\
$r_i\in\{\ATCO,\PILOT\}$ & 화자 역할 & $\hat c_i$ & 콜사인 언급(문자열, 생략 시 $\varepsilon$)\\
$\iota=(\kappa,v,\mu,q,g)$ & 지시 개념(유형·값·단위·한정어·대상) & $I_i$ & 송신 $i$의 지시 집합\\
$f_i \in V_{\mathrm{Flight}}$ & 해석된 항공편 노드 & $\Tcmd$ & 지시 유형 집합\\
$\Treq\subseteq\Tcmd$ & 복창 필수 유형 & $\Tst\subseteq\Tcmd$ & 상태형(유효 구간을 가지는) 유형\\
$\Tadv$ & 정보·권고(advisory) 유형 & $\G_{<t}$ & 시각 $t$ 이전까지 누적된 그래프\\
$\G|_t$ & 시각 $t$의 스냅샷 & $\Delta_{\mathrm{rb}}$ & 복창 허용 시간창\\
$\KR$ & 규칙 베이스 (\texttt{rules.json}) & $\Cfg$ & 구축 설정 (\texttt{kg\_build.json})\\
\bottomrule
\end{tabular}
\end{center}

%==========================================================
\section{데이터 명세}\label{sec:data}
%==========================================================
\subsection{구축용 데이터 $D_{\mathrm{US}}$: HF LiveATC 원본 아카이브}
\begin{itemize}
\item \textbf{출처}: Hugging Face \texttt{nyuuzyou/liveatc} \cite{hfliveatc}. LiveATC.net의 30분 단위 MP3 아카이브이며, 2024-08-26 하루분이다. 수집 한도 때문에 일부만 들어 있다.
\item \textbf{규모·구성 (직접 확인)}: 213개 공항 아카이브, 45.5\,GB. \textbf{전사가 없는 음성 전용} 데이터라 M1a--M1b의 STT 단계가 필수다.
\item \textbf{국가 구성과 공항 범위}: 미국·캐나다·유럽 공항이 섞여 있다(내용 있는 파일 20,742개 중 미국 117개 피드 10,029개 48\%, 캐나다 31.5\%, 유럽 9개국 19.2\%, 레위니옹·마다가스카르 0.9\%). 미국 공항은 코드 알파벳 앞쪽 일부만 들어 있어 KSFO, KJFK 같은 주요 공항이 없고, 일본(\texttt{RJ*})·한국(\texttt{RK*}) 공항도 없다. 그래서 구축 대상 공항을 KSFO에서 \textbf{KATL(애틀랜타)}로 바꾼다.
\item \textbf{KATL 피드}: \cmd{Gnd-All}, \cmd{Del}, \cmd{Dep-North/South}, \cmd{App-Final}(3개), \cmd{ATIS-Arr/Dep}, \cmd{Company}. \textbf{Tower 피드는 없다}.
\item \textbf{결손}: \cmd{Gnd-All}은 48개 파일 중 \textbf{11개만 실제 음성}이다(00:00--05:30Z, 약 5.5시간). 나머지는 0바이트다. 다른 피드는 24시간이 모두 있다.
\item \textbf{샘플 확인}: \cmd{Gnd-All} 02:30Z 파일의 5분을 전사했다.
  \begin{itemize}
  \item 디지털 무음 0\%, 음성 에너지 대역 320--1240\,Hz (16\,kbps).
  \item taxi via / hold short / give way / monitor tower 같은 \textbf{지시$\to$복창 쌍이 5분에 10쌍 이상} 있었다.
  \item 오인식도 그대로 관찰되었다: \texttt{line of ground}(Atlanta Ground), \texttt{Fox, shot}(Foxtrot), \texttt{19-1}(119.1), \texttt{Alpha 265}$\leftrightarrow$\texttt{Delta 265}.
  \end{itemize}
\end{itemize}

\paragraph{문제정의와의 적합성.}
\begin{center}\small
\begin{tabularx}{\textwidth}{@{}>{\raggedright\arraybackslash}p{4.2cm}>{\raggedright\arraybackslash}X@{}}
\toprule
문제정의 요소 & HF LiveATC(KATL)에서의 대응\\\midrule
음성 전용 입력, STT 필요 (\ch{C1}, \ch{C2}) & 적합. 전사가 없고, 협대역·도메인 오인식이 실제로 관찰됨\\
지시–복창 루프 (RQ1, RQ3) & 적합. \cmd{Gnd-All}에 taxi/hold short/give way 지시와 복창이 밀집\\
콜사인 오인식 (\ch{C4}) & 적합. \texttt{Alpha}/\texttt{Delta} 혼동이 재현됨\\
맥락 $\Ctx$ (\ch{C9}) & \textbf{개선}. 파일명에서 공항·시설·UTC 시각을 자동으로 얻음 ($\tau_0$가 절대 시각)\\
시간 KG (\ch{C10}) & \textbf{개선}. 편집된 유튜브 샘플과 달리 30분 연속 녹음이라 지시$\to$복창$\to$이관 흐름이 끊기지 않음. 여러 피드를 UTC로 맞춰 \cmd{Del}$\to$\cmd{Gnd}$\to$\cmd{Dep} 시설 간 이관도 연결 가능\\
착륙 허가·활주로 중복 ($\rho_4$, $q_1$) & \textbf{부분 부적합}. KATL Tower 피드가 없음. 대체: \cmd{CROSS\_RUNWAY}/\cmd{HOLD\_SHORT} 기반 활주로 질의로 바꾸거나, Tower가 있는 공항(예: \cmd{KABQ-Twr-App}, \cmd{KBDL-Twr})을 보조로 사용\\
\bottomrule
\end{tabularx}
\end{center}

\paragraph{필터링 규칙.} 사용할 시설을 $\Phi_{\mathrm{use}}=\{\cmd{Gnd},\cmd{Del},\cmd{Dep},\cmd{App},\cmd{Twr}\}$로 둔다.
\begin{itemize}
\item \cmd{CTAF}: 관제사가 없는 비관제 공항 주파수라 지시–복창 개념이 성립하지 않으므로 제외한다.
\item \cmd{Company}: 항공사 운항 통신이라 제외한다.
\item \cmd{ATIS}: 단일 화자 방송이므로 송신으로는 쓰지 않는다. 대신 사용 활주로·정보 부호(예: \texttt{information Papa})를 $\Ctx$에 넣는 \emph{선택적 맥락원}으로 쓴다.
\end{itemize}

\paragraph{이용 조건.} LiveATC 약관상 \textbf{개인·비상업 용도}이며, 재배포에는 허가가 필요하다 \cite{hfliveatc}. 수업 내 분석·데모에는 쓰되, 음성 원본과 전사 전문을 공개 저장소에 올리지 않는다.

\subsection{후보 데이터 비교와 선정 근거}\label{sec:datacmp}
\paragraph{문제정의에서 도출한 데이터 요구사항.} 데이터는 연구 질문과 기술적 문제를 \emph{실제로 재현}할 수 있어야 한다.
\begin{itemize}
\item[\textbf{R1}] \textbf{미국(FAA) 실제 운영 교신}: 구축 범위(RQ1--RQ3)와 FAA 복창 규정 \cite{faa7110}, 미국 특수성(\ch{C11}) 때문.
\item[\textbf{R2}] \textbf{관제사와 조종사 양측 음성}: 지시--복창 루프(RQ3, \ch{C6}) 자체가 양측 발화를 요구.
\item[\textbf{R3}] \textbf{시간 연속성}: 지시$\to$복창$\to$이관이 끊기지 않는 연속 녹음과 시각 정보(\ch{C8}, \ch{C10}).
\item[\textbf{R4}] \textbf{맥락 메타데이터}: 공항·주파수·절대 시각(\ch{C9}). 감시 데이터(ADS-B)가 있으면 더 좋음(RQ2, \cite{nigmatulina2022}).
\item[\textbf{R5}] \textbf{즉시·무료 확보}: 학기 일정 안에 바로 쓸 수 있어야 함.
\item[\textbf{R6}] \textbf{RQ4 연계}: 비교 지역(일본) 데이터와 출처가 같으면 교란 요인이 줄어듦.
\item[\textbf{R7}] (선택) \textbf{정답 전사}: 진단(\md{M7})과 점검에 유용하나, 소량은 직접 달 수 있으므로 필수는 아님.
\end{itemize}

\begin{center}\footnotesize
\setlength{\tabcolsep}{3pt}
\begin{tabularx}{\textwidth}{@{}>{\raggedright\arraybackslash}p{2.6cm}>{\raggedright\arraybackslash}p{3.4cm}ccccccc>{\raggedright\arraybackslash}X@{}}
\toprule
데이터 & 내용 & R1 & R2 & R3 & R4 & R5 & R6 & R7 & 판단\\\midrule
\textbf{HF LiveATC} \cite{hfliveatc} & 미국·캐나다·유럽 혼합(미국 117개 피드, 48\%), 30분 연속 MP3, 2024-08-26 & $\bullet$ & $\bullet$ & $\bullet$ & $\circ$ & $\bullet$ & $\bullet$ & -- & \textbf{주 구축 데이터}\\
LiveATC 아카이브 직접 수집 & 임의 공항(예: KSFO Tower), 최근 기간만 보관 & $\bullet$ & $\bullet$ & $\bullet$ & $\circ$ & $\circ$ & $\bullet$ & -- & 보조 (Tower 보강, 일본 검증)\\
TartanAviation \cite{tartan2025} & KAGC Tower(유관제)·KBTP CTAF, 3{,}374\,h, ADS-B 동시 수집 & $\bullet$ & $\bullet$ & $\circ$ & $\bullet$ & $\circ$ & -- & -- & 보조 후보 (Tower + ADS-B). 단 소형기(GA) 공항\\
LDC ATC0 \cite{atc0} & DFW/BOS/DCA, 약 70\,h, 시각 포함 전사 & $\bullet$ & $\bullet$ & $\bullet$ & $\circ$ & -- & -- & $\bullet$ & 학교 LDC 회원이면 정답 전사 보조. 1990년대·유료\\
ATCO2 \cite{atco2} & 유럽 위주 실제 교신. 정답 4\,h(1\,h 무료)+자동전사 5{,}281\,h, NER 주석 & -- & $\bullet$ & $\circ$ & $\circ$ & $\circ$ & -- & $\bullet$ & 구축 부적합(R1). 추출 스키마 참고용\\
UWB-ATCC \cite{uwbatcc} & 체코 실제 교신, 약 13\,h, 전사 & -- & $\bullet$ & -- & -- & $\bullet$ & -- & $\bullet$ & 부적합 (R1, R3)\\
ATCOSIM \cite{atcosim} & 유럽 시뮬레이터, 약 10\,h, \textbf{관제사 음성만} & -- & -- & -- & -- & $\bullet$ & -- & $\bullet$ & 부적합. 복창이 없음(R2)\\
Airbus ATC \cite{airbusatc} & 프랑스 억양 영어, 약 59\,h & -- & $\bullet$ & $\circ$ & -- & -- & -- & $\bullet$ & 부적합 (비공개)\\
ATCSpeech \cite{atcspeech} & 중국 실제 교신, 중국어+영어 혼용 & -- & $\bullet$ & $\circ$ & -- & -- & -- & $\bullet$ & 구축 부적합. 비영어권 확장(향후 과제) 참고\\
YouTube-ATC \cite{youtubeatc} & VATSIM/IVAO 스트림 2{,}000\,h+, 17개국, 자동 전사 & -- & $\bullet$ & $\bullet$ & -- & $\circ$ & $\circ$ & $\circ$ & 구축 부적합(모의). \textbf{KR-1안(VATSIM) 사용의 선례}\\
팀 유튜브 샘플 & 편집된 모음, 248\,s & $\bullet$ & $\bullet$ & -- & -- & $\bullet$ & -- & -- & 문제 관찰용으로만 사용 (\S\ref{sec:challenges})\\
\bottomrule
\end{tabularx}
\end{center}
{\footnotesize $\bullet$ 충족, $\circ$ 부분 충족, -- 미충족. ``R4 $\circ$''는 공항·주파수·시각은 있으나 감시 데이터가 없다는 뜻이다.}

\paragraph{선정 논리.}
\begin{enumerate}
\item \textbf{R1--R5를 동시에 만족하는 것은 LiveATC 계열뿐이다.} 정답 전사가 있는 공개 코퍼스는 대부분 유럽 데이터이거나(ATCO2, UWB-ATCC), 시뮬레이터 데이터이거나(ATCOSIM), 비공개다(Airbus). 미국 데이터인 LDC ATC0는 유료이고 30년 전 녹음이다.
\item \textbf{전사가 없다는 점은 약점이 아니라 문제정의와 일치하는 조건이다.} 문제정의의 입력은 ``음성으로만 존재하는 교신''이다. STT 오류(\ch{C1}, \ch{C2}, \ch{C4})를 KG로 탐지·교정하는 것이 RQ2이므로, 깨끗한 정답 전사로 시작하면 오히려 RQ2를 검증할 수 없다.
\item \textbf{공개 수신 음성은 ATC 연구의 표준 원천이다.} ATCO2는 공개 VHF 수신 음성을 대규모로 모아 코퍼스를 만들었고 \cite{atco2}, 맥락 기반 ATC 음성 인식 연구는 LiveATC에서 모은 테스트셋을 썼다 \cite{zuluaga2021}. 우리 데이터 선택은 이 선행 관행과 같은 방향이다.
\item \textbf{연속 녹음이라 Temporal KG에 맞다.} 발화 단위로 잘린 코퍼스(UWB-ATCC 등)는 발화 간 관계(\ch{C8})와 유효 구간(\ch{C10})을 복원할 수 없다.
\item \textbf{RQ4와 출처가 같다.} 일본 검증 데이터도 LiveATC에서 가져오면 수집 장비·형식 차이가 아닌 \emph{지역 차이}만 남는다(R6).
\end{enumerate}

\paragraph{남는 한계와 보완.}
\begin{itemize}
\item \textbf{감시 데이터 없음 (R4 부분)}: OpenSky 과거 데이터로 같은 시각의 KATL 주변 항적을 붙일 수 있다(연구용 접근 신청 필요, 선택). 붙이지 않으면 $\mathcal{F}_{\mathrm{act}}$는 그래프 내부 콜사인으로만 구성한다(\md{M1b}, \md{M3}).
\item \textbf{Tower 부재}: 착륙 허가 시나리오가 필요하면 TartanAviation KAGC Tower(ADS-B 포함) 또는 LiveATC 아카이브의 Tower 피드를 보조로 쓴다. KAGC는 소형기 공항이라 콜사인 형태(\texttt{November 123AB})와 교통 유형이 KATL과 다르다는 점을 명시해야 한다.
\item \textbf{정답 없음 (R7)}: \md{M7} 진단용으로 국가별 15--20개 송신에 직접 정답을 단다. 학교가 LDC 회원이면 ATC0 전사를 정답 비교용으로 일부 쓸 수 있다.
\end{itemize}

\subsection{검증용 데이터 (RQ4)}
\begin{center}\small
\begin{tabularx}{\textwidth}{@{}l>{\raggedright\arraybackslash}p{3.5cm}>{\raggedright\arraybackslash}X@{}}
\toprule
지역 & 출처 & 문제정의 관점의 주의점\\\midrule
JP & LiveATC 일본 공항 피드 (HF 데이터셋에는 없음. LiveATC 아카이브에서 직접 수집) & 실제 운영 교신이고 ICAO 용어를 쓴다. 미국 데이터와 같은 출처라 \emph{지역 차이만} 비교할 수 있다. 피드 존재 여부와 보관 기간은 확인 필요\\
KR-1안 & VATSIM 한국 관제 연습 교신 & 보안 문제는 없다. VATSIM/IVAO 음성은 ATC 음성 인식 연구에 이미 쓰인 선례가 있다 \cite{youtubeatc}. 그러나 \textbf{시뮬레이션}이라 지역 차이 외에 \emph{실제 대 모의} 차이가 섞인다(교란 요인). 다른 참가자 음성을 녹음하므로 VATSIM 규정과 참가자 동의 확인이 필요\\
KR-2안 & 국내 실제 녹취 (보안 자료) & 실제 운영 교신이라 비교에 가장 적합하다. 그러나 \textbf{외부 LLM API로 전송 금지}, 원본·전사 제출 금지, 집계 지표만 보고. 소유 기관 허가가 필요\\
\bottomrule
\end{tabularx}
\end{center}
권장: KR-1안을 기본으로 하고, 2안은 허가가 나면 \emph{로컬에서만} 처리하는 보조 조건으로 둔다. 1안의 교란 요인을 분리하려면 같은 VATSIM에서 미국 공역 교신 소량을 대조군으로 쓸 수 있다(선택).

%==========================================================
\section{기술적 문제 정의 (Challenges)}\label{sec:challenges}
%==========================================================
아래 문제들은 문제정의 문서의 논증과, 샘플 음성을 Whisper large-v3로 전사·측정한 결과에서 도출했다.
각 문제는 \S\ref{sec:modules}의 모듈과 대응된다(\S\ref{sec:trace}).

\begin{center}\small
\begin{tabularx}{\textwidth}{@{}l>{\raggedright\arraybackslash}p{3.6cm}>{\raggedright\arraybackslash}X l@{}}
\toprule
ID & 기술적 문제 & 증거 (샘플 측정 / 선행연구) & 대응 RQ\\\midrule
\ch{C1} & 무음 구간 ASR 환각 & 248초 중 \textbf{57.5\%가 디지털 무음}. VAD 없이 전사하면 \texttt{"Thanks for watching."}, \texttt{"you"}가 생성됨. Whisper가 무음·긴 휴지에서 환각을 일으킨다는 보고 \cite{koenecke2024} & RQ1\\
\ch{C2} & 협대역 음성과 도메인 어휘 오인식 & 에너지 90\%가 240--1720\,Hz에 집중. \texttt{Terra}(Tower), \texttt{north-south}(NorCal), \texttt{skip away}(give way). 일반 Whisper는 ATC에서 성능이 낮음 \cite{atco2,vandoorn} & RQ2\\
\ch{C3} & 화자 경계·역할 불명 & 송신 간격 0.1--0.5\,s. 관제사와 조종사 발화가 한 세그먼트로 병합됨. 텍스트 기반 역할 판별 \cite{bertraffic} & RQ1\\
\ch{C4} & 콜사인 오인식·생략 & 같은 음성이 \texttt{Delta 2398} / \texttt{Alpha 2398}로 다르게 전사됨. 콜사인 없는 복창 \texttt{"Alpha, alpha one, two."}. 콜사인 문제는 통신 사고의 주요 기여 요인 \cite{eurocontrol,cardosi94} & RQ2, RQ3\\
\ch{C5} & 숫자·표기 변이 & \texttt{two two right} / \texttt{22 right}, \texttt{20-69} / \texttt{2069} & RQ2\\
\ch{C6} & 바꿔 말하는 복창 (단어 비교로는 실패) & 복창은 순서·어휘·콜사인 위치가 지시와 다름. 단어 단위 비교는 비효과적 \cite{helmke2021} & RQ3\\
\ch{C7} & 지시와 정보(advisory)의 구분 & \texttt{heavy departed}, \texttt{parallel runway}는 복창 대상이 아님 & RQ3\\
\ch{C8} & 발화 간 관계 (지시→복창→이관) & 의미는 여러 송신에 걸쳐 완성됨. 발화 단위 주석만으로는 표현 불가 \cite{kleinert2021,helmke2018} & RQ1\\
\ch{C9} & 맥락 전환과 정적 지식 필요 & 22R/Juliet(타 공항)과 SFO 28L/28R이 한 클립에 혼재. 맥락 정보가 검증에 필요 \cite{helmke2021,nigmatulina2022} & RQ2\\
\ch{C10} & 사실의 시간 유효성 & 착륙 허가 등은 구간 동안만 유효. 정적 KG는 과거와 현재 허가를 구분하지 못함 \cite{cai2024tkg,knez2023} & RQ1, RQ3\\
\ch{C11} & 라벨 부족과 미국·유럽 차이 & 미국 교신에 대한 의미 주석 데이터가 없음. 미국과 유럽의 어휘·개념 차이 \cite{chen2023} & RQ1, RQ2\\
\ch{C12} & 오류 전파와 설명 가능성 & ASR·LLM 오류가 KG로 전파됨. 안전 분야는 판정 근거 제시가 필요함 \cite{pan2024} & RQ3\\
\ch{C13} & LLM 출력의 형식 오류와 환각 & 자유 생성은 스키마 위반이나 입력에 없는 개체를 만들 수 있음 \cite{willard2023,itext2kg} & RQ1\\
\ch{C14} & 지역 전이 (비영어권 억양, ICAO와 FAA 용어 차이) & 지역 데이터로 만든 모델은 다른 지역에서 잘 작동하지 않음 \cite{chen2023,usatc2025}. QNH/altimeter, holding point/hold short 같은 용어 차이. 국내선 자국어 혼용 가능성(샘플 확인 필요) & RQ4\\
\ch{C15} & 데이터 이용 제약 & LiveATC 약관(비상업·재배포 금지) \cite{hfliveatc}, 한국 실녹취의 보안 등급. 외부 API로 음성·전사를 보내면 반출에 해당함 & RQ4\\
\bottomrule
\end{tabularx}
\end{center}

%==========================================================
\section{설계 결정과 적합성 근거}\label{sec:design}
%==========================================================
각 결정은 ``대안 $\to$ 선택 $\to$ 우리 문제에 맞는 이유'' 순서로 정리한다.

\begin{center}\small
\begin{tabularx}{\textwidth}{@{}l>{\raggedright\arraybackslash}p{3.1cm}>{\raggedright\arraybackslash}p{3.0cm}>{\raggedright\arraybackslash}X@{}}
\toprule
& 선택 & 기각한 대안 & 근거 (문제 ID와 문헌)\\\midrule
D1 & \textbf{계단식 파이프라인} (ASR $\to$ 의미추출 $\to$ KG) & 음성$\to$KG 종단간(end-to-end) 모델 & 학습 데이터가 없고(\ch{C11}) 중간 산출물로 오류를 추적해야 함(\ch{C12}). DLR·ATCO2의 ATC 음성 이해 연구도 ASR$\to$NLU 계단식 구조를 씀 \cite{kleinert2021,atco2}\\
D2 & \textbf{스키마 기반(온톨로지 기반) 추출} & 스키마 없는 개방형 추출 (EDC 개방 단계, GraphRAG) & 복창 비교(\ch{C6})에는 \emph{정규화된 개념}이 필요함. 검증된 도메인 온톨로지와 공식 사양이 이미 있음 \cite{dlr2025,helmke2018,chen2023,keller2017}. 온톨로지 기반 생성은 일관성과 해석 가능성을 높임 \cite{feng2024}. LLM-KGC 서베이의 schema-based 패러다임에 해당 \cite{llmkgc2025}. GraphRAG는 질의 요약용 자유 개체 그래프라 명령 의미를 담지 못함 \cite{graphrag}\\
D3 & \textbf{LLM few-shot 추출} & BERT NER 미세조정 & 미국 주석 데이터가 없음(\ch{C11}). ATCO2 NER은 1시간 학습으로 F1 0.66 \cite{atco2}. LLM 추출은 대량 주석 없이 micro-F1 0.95(사람 전사)/0.86(Whisper 전사) \cite{mekerishvili2025}\\
D4 & \textbf{JSON Schema 제약 디코딩} + 사후 검증 & 자유 생성 후 파싱 & 출력이 문법상 항상 스키마를 만족하도록 강제(\ch{C13}) \cite{willard2023}. \emph{[v2 구현: 사후 검증 + 어휘 단서 + 규칙 보충. \S\ref{sec:impl}]}\\
D5 & \textbf{증분 구축 + 개체 해소} & 문서 단위 일괄 구축 후 병합 & 송신은 시간순 스트림임. 증분 개체·관계 해소로 중복과 환각 개체를 억제 \cite{itext2kg}(\ch{C4}, \ch{C13})\\
D6 & \textbf{사건 노드 reification} & 평면 트리플 $(s,p,o)$ & 지시는 n-항 관계이고(\ch{C8}) 사건 중심 TKG가 텍스트 기반 시간 KG 구축의 주류 \cite{naryw3c,knez2023}\\
D7 & \textbf{이중 시간(bi-temporal) + 삭제 대신 무효화} & 스냅샷 KG, TKG 임베딩(TTransE 류) & 사실의 유효 구간(\ch{C10})과 관측 시각을 분리. 갱신 시 이력을 보존 \cite{zep2025,atom2025}. 임베딩 기반 TKG는 예측·완성용이며 대량 데이터와 정량 평가가 필요해 우리 과제(탐지)와 맞지 않음 \cite{cai2024tkg}\\
D8 & \textbf{기호 규칙 기반 탐지} & LLM 복창 분류기 (SCOPE) & 규칙은 규정(FAA JO 7110.65)에서 직접 도출 가능하고, 근거 경로를 제시할 수 있으며(\ch{C12}), 학습 데이터가 필요 없음(\ch{C11}). LLM 분류기는 항공편 이력이나 정적 KB를 쓰지 못함 \cite{scope2026,faa7110}\\
D9 & \textbf{맥락 어휘 프롬프트 바이어싱} & 미국 ATC로 Whisper 미세조정 & 라벨이 없고 시간이 부족함(\ch{C11}). 감시 데이터 기반 콜사인 바이어싱만으로 콜사인 인식이 크게 개선됨 \cite{nigmatulina2022}. 미세조정은 선택적 비교 실험으로 둠 \cite{vandoorn}\\
D10 & \textbf{속성 그래프(Neo4j)} + 온톨로지 IRI 매핑 & RDF 트리플스토어 & 데모용 시간 질의(Cypher)가 쉽고, iText2KG·Graphiti 같은 최신 구축 도구가 Neo4j를 지원함 \cite{itext2kg,zep2025}. 의미 정합성은 노드 라벨을 atmonto/SESAR 개념 IRI에 매핑해 유지. 제약 $\Sigma$는 SHACL 형태로 정의해 검사 \cite{shacl}. \emph{[v2 구현: JSON 그래프 + Neo4j 내보내기. \S\ref{sec:impl}]}\\
\bottomrule
\end{tabularx}
\end{center}

\paragraph{프레임워크 정합성.}
제안 구조는 두 가지 기존 틀과 대응된다.
\begin{itemize}
\item \textbf{LLM-KGC 서베이의 3단계} \cite{llmkgc2025}: 온톨로지 공학 = \md{M0}, 지식 추출 = \md{M2b}, 지식 융합 = \md{M3}--\md{M4}.
\item \textbf{Endsley의 상황 인식 3단계} \cite{endsley1995}: 지각(Level 1) = \md{M1a}--\md{M2b}, 이해(Level 2) = \md{M3}--\md{M5}, 예측(Level 3, 제한적) = \md{M6}.
\end{itemize}
Zep/Graphiti의 \emph{episodic--semantic 부분그래프} 구분 \cite{zep2025}과도 대응된다. 송신(\cmd{Transmission})은 episode, 항공편·지시는 semantic 노드에 해당한다.

%==========================================================
\section{시스템 아키텍처}
%==========================================================
\begin{figure}[h!]
\centering
\resizebox{\textwidth}{!}{%
\begin{tikzpicture}[
  font=\small,
  mod/.style={draw, rounded corners=3pt, thick, fill=white, minimum width=2.55cm, minimum height=1.25cm, align=center},
  io/.style={draw, thick, fill=gray!12, minimum width=2.1cm, minimum height=1.25cm, align=center},
  resb/.style={draw, thick, dashed, rounded corners=3pt, fill=res, minimum width=2.6cm, minimum height=3.6cm, align=center},
  arr/.style={-{Stealth[length=2.6mm]}, thick},
  lab/.style={font=\footnotesize, fill=white, inner sep=1pt},
  chl/.style={font=\scriptsize\color{red!60!black}}
]
% row 1
\node[io]  (in) at (0,0) {\textbf{입력}\\$a,\;\Ctx$};
\node[mod] (m1) at (3.3,0) {\textbf{M1a}\\송신 분할};
\node[mod] (m2) at (6.6,0) {\textbf{M1b}\\맥락 바이어싱\\ASR + 환각필터};
\node[mod] (m3) at (9.9,0) {\textbf{M2a}\\화자 역할\\판정};
\node[mod] (m4) at (13.2,0) {\textbf{M2b}\\온톨로지 기반\\LLM 의미추출};
% row 2
\node[mod] (m5) at (13.2,-3.4) {\textbf{M3}\\정규화·grounding\\·개체 해소};
\node[mod] (m6) at (9.9,-3.4) {\textbf{M4}\\증분 Temporal\\KG 구축};
\node[mod] (m7) at (6.6,-3.4) {\textbf{M5}\\규칙 추론·\\이상 탐지};
\node[mod] (m8) at (3.3,-3.4) {\textbf{M6}\\시점 질의·\\시각화};
\node[io]  (out) at (0,-3.4) {\textbf{출력}\\$\G,\;\Yset,\;Q$};
\node[mod, dashed, fill=gray!8] (m9) at (3.3,-7.6) {\textbf{M7}\\진단 하네스\\(RQ4)};
\node[io] (out2) at (0,-7.6) {\textbf{RQ4 진단}\\$\mathrm{CSR},\mathrm{CmdF1},m^\dagger$};
% resources
\node[resb] (m0) at (17.9,-1.7) {\textbf{M0}\\정적 자원\\(오프라인)\\[3pt]$\Ont=(\mathcal{T},\mathcal{R},\Sigma)$\\$\KS=(V_S,E_S)$\\$\KR,\ \Cfg,\ \Vdom$};

% main arrows
\draw[arr] (in) -- node[lab,above=2pt]{$a,\Ctx$} (m1);
\draw[arr] (m1) -- node[lab,above=2pt]{$\{(s_i,x_i)\}$} (m2);
\draw[arr] (m2) -- node[lab,above=2pt]{$(w_i,c_i)$} (m3);
\draw[arr] (m3) -- node[lab,above=2pt]{$(w_i,r_i)$} (m4);
\draw[arr] ($(m4.south)+(0.95,0)$) -- node[lab,right=2pt]{$u_i$} ($(m5.north)+(0.95,0)$);
\draw[arr] (m5) -- node[lab,above=2pt]{$\tilde u_i$} (m6);
\draw[arr] (m6) -- node[lab,above=2pt]{$\G_{t}$} (m7);
\draw[arr] (m7) -- node[lab,above=2pt]{$\G_t,\Yset$} (m8);
\draw[arr] (m8) -- node[lab,above=2pt]{} (out);
% feedback
\draw[arr, dashed, blue!60!black] ($(m6.south)+(0.7,0)$) -- ++(0,-0.75) -| node[lab,pos=0.75,below]{$\G_{<t_i}$} ($(m5.south)+(-0.7,0)$);
% resources arrows
\draw[arr, violet] (m0.north) |- ++(0,1.75) -| node[lab,pos=0.25,above]{$\Vdom$, $\mathcal{F}_{\mathrm{act}}$ (콜사인·시설·유도로 어휘)} (m2.north);
\draw[arr, violet] (m0.west |- m4) -- node[lab,above]{$\Ont$, 예시 $\mathcal{E}$} (m4.east);
\draw[arr, violet] (m0.west |- m5) -- node[lab,above]{$\KS$} (m5.east);
\draw[arr, violet] (m0.south) |- ++(0,-2.05) -| node[lab,pos=0.25,below]{$\Sigma$, $\Tst$, $\KR$ (rules.json), $\Cfg$ (kg\_build.json)} ($(m6.south)+(-0.7,0)$);

\draw[arr, dashed, gray] (m8.south) -- node[lab,right=2pt,pos=0.72]{중간 산출물 로그 $\mathcal{J}$ (M1a--M5)} (m9.north);
\draw[arr, dashed, gray] (m9) -- (out2);
\node[chl, anchor=west] at ($(m9.east)+(0.1,0)$) {C14, C15};
% challenge annotations
\node[chl, below=1pt of m1] {C1, C9};
\node[chl, below=1pt of m2] {C1, C2, C11, C14};
\node[chl, below=1pt of m3] {C3};
\node[chl, below=1pt of m4, xshift=-0.35cm] {C6, C7, C11, C13, C15};
\node[chl, above=1pt of m5, xshift=-0.35cm] {C2, C4, C5, C9};
\node[chl, above=1pt of m6] {C8, C10, C12, C13};
\node[chl, above=1pt of m7] {C4, C6, C7, C12};
\node[chl, above=1pt of m8] {C10, C12};

% SA background bands
\begin{scope}[on background layer]
\node[fill=perc, rounded corners, fit=(m1)(m4), inner xsep=8pt, inner ysep=20pt, label={[font=\footnotesize\bfseries,anchor=south west]north west:Level 1 지각 (Perception)}] {};
\node[fill=comp, rounded corners, fit=(m5)(m7), inner xsep=8pt, inner ysep=22pt, label={[font=\footnotesize\bfseries,anchor=north west]south west:Level 2 이해 (Comprehension)}] {};
\node[fill=proj, rounded corners, fit=(m8), inner xsep=6pt, inner ysep=22pt, label={[font=\footnotesize\bfseries]below:Level 3 예측/질의}] {};
\end{scope}
\end{tikzpicture}}
\caption{시스템 모듈 구성. 실선은 데이터 흐름, 점선(파랑)은 누적 그래프를 개체 해소의 맥락으로 되먹이는 경로, 보라색은 정적 자원 \md{M0}의 공급 경로다. 각 모듈 아래·위의 빨간 표기는 그 모듈이 해결하는 기술적 문제(\S\ref{sec:challenges})이고, 배경색은 Endsley의 상황 인식 단계 \cite{endsley1995}다.}
\label{fig:arch}
\end{figure}

전체 처리는 송신 단위 스트림으로 진행된다. $i=1,\dots,n$에 대해
\begin{equation}
\G_{t_i} \;=\; \mathrm{M4}\Big(\G_{t_{i-1}},\; \mathrm{M3}\big(\mathrm{M2b}(\mathrm{M2a}(\mathrm{M1b}(x_i))),\; \KS,\; \G_{t_{i-1}}\big)\Big),
\qquad \G_{t_0} = \KS,
\end{equation}
이고, 이상 집합은 $\Yset = \mathrm{M5}(\G_{t_n})$이다.

%==========================================================
\section{모듈 정의}\label{sec:modules}
%==========================================================
각 모듈은 \textbf{[해결 문제] / [입력] / [출력] / [정의] / [방법·근거]} 순서로 기술한다.

%----------------------------------------------------------
\subsection{M0. 정적 자원 구축 (Ontology \& Static KB) --- 오프라인}
\textbf{[해결 문제]} \ch{C9} 맥락·정적 지식, \ch{C11} 미국판 개념 체계, \ch{C2} 도메인 어휘 공급.

\textbf{[입력]}
\begin{itemize}
\item $D_{\mathrm{apt}}$: 공항 정적 정보. OurAirports의 활주로와 OpenStreetMap의 유도로 이름 \textbf{목록}(좌표·연결 관계는 쓰지 않음).
\item $D_{\mathrm{tel}}$: 항공사 telephony$\leftrightarrow$ICAO 표. FAA JO 7340.2 \cite{faa7340}.
\item $D_{\mathrm{fac}}$: 관제시설·주파수 목록 (OurAirports. 예: KATL \cmd{GND} 121.65, \cmd{TWR} 119.1).
\item $D_{\mathrm{onto}}$: \textbf{DLR 온톨로지 공식 사양 v3.00} \cite{dlr2025}. 설명서 PDF(지시 구조 Figure 2-1, JSON 키 5.1장)와 명령 표 PDF(16개 그룹, Type (Main) 91개, 126줄과 매개변수 형식).\footnote{설명서 9쪽은 Cleared/Cancel을 합쳐 15개 그룹으로 나열한다.} 배경 문헌: SESAR 논문 \cite{helmke2018,kleinert2021}, MITRE 비교 \cite{chen2023}, NASA atmonto \cite{keller2017}.
\item $D_{\mathrm{reg}}$: \textbf{FAA Order JO 7110.65BB} \cite{faa7110}. 명령별 표준 문구(PHRASEOLOGY)와 복창 규정 §2-4-3. KR/JP의 ICAO 규정은 RQ4 단계에서 추가한다(현재 미검증).
\item 지역 매개변수 $\rho\in\{\mathrm{US},\mathrm{KR},\mathrm{JP}\}$: 공항, 복창 규정, 용어 표면형(예: ICAO \texttt{holding point}, \texttt{QNH} $\leftrightarrow$ FAA \texttt{hold short}, \texttt{altimeter}).
\end{itemize}

\textbf{[출력]}
\begin{align}
\Ont &= (\mathcal{T},\mathcal{R},\Sigma),\quad \mathcal{T} = \mathcal{T}_{\mathrm{ent}} \cup \Tcmd,\\
\KS &= (V_S, E_S),\quad V_S \subseteq \{\cmd{Airport},\cmd{Runway},\cmd{Taxiway},\cmd{Facility},\cmd{Airline}\}\text{-유형 노드},\\
\Vdom(\alpha_{\mathrm{apt}}) &= \{\mathrm{tel}(l) : l\in V_{\cmd{Airline}}\} \cup \{\mathrm{name}(v): v \in V_S,\ v \sqsubseteq \alpha_{\mathrm{apt}}\},\\
\KR &= \big(\{\Treq^{\rho}\}_{\rho},\ \mathrm{sup},\ \mathcal{T}_{\mathrm{rwy}}\big),\qquad \Cfg = \big(\Delta_{\mathrm{rb}},\ T_{\max}\big).
\end{align}
구현 파일은 다음과 같다. $\Ont$와 $\KR$은 \texttt{kb/build\_ontology.py}가 PDF 원문과 대조해 생성하고, $\KS$는 \texttt{kb/build\_static\_kb.py}가 생성하며, $\Cfg$와 항공사 표는 사람이 작성한다.
\begin{center}\small
\begin{tabularx}{\textwidth}{@{}l l>{\raggedright\arraybackslash}X l@{}}
\toprule
기호 & 파일 & 내용 & 근거\\\midrule
$\Ont$ & \texttt{configs/ontology.json} & 개체·관계 유형, 명령 20종(DLR 표의 행·값 형식, FAA 표준 문구), DLR 지시 구조와 16개 그룹 목록, 오탐 방지 어휘 단서 & DLR v3.00, FAA\\
$\KR$ & \texttt{configs/rules.json} & 지역 선택, 복창 등급과 필수 목록(FAA 13종, ICAO 12종), 지시 대체 관계, 활주로 점유 유형 & FAA §2-4-3, 지시 의미\\
$\Cfg$ & \texttt{configs/kg\_build.json} & $\Delta_{\mathrm{rb}}=20$\,s, $T_{\max}(\kappa)$ 180--1800\,s (기본 600\,s) & 설계값(KATL 실측 참고)\\
$\KS$ & \texttt{data/static\_kb/katl.json} & 활주로 5개(이름 10개), 유도로 102개, 관제시설 12개 & OurAirports, OSM\\
$D_{\mathrm{tel}}$ & \texttt{configs/airlines.json} & 호출 이름 $\to$ ICAO 코드 44개 & FAA JO 7340.2\\
\bottomrule
\end{tabularx}
\end{center}

\textbf{[정의]} 지시 유형 $\Tcmd$는 DLR 온톨로지의 명령 유형을 기본으로 하고, 미국에만 있는 어휘(\texttt{niner}, \texttt{altimeter} 등)를 표면형 사전에 추가한다 \cite{chen2023}.
\begin{align*}
\Tcmd = \{&\cmd{CLEARED\_TO\_LAND}, \cmd{CLEARED\_FOR\_TAKEOFF}, \cmd{LINE\_UP\_AND\_WAIT}, \cmd{TAXI\_VIA}, \cmd{CROSS\_RUNWAY},\\
&\cmd{HOLD\_SHORT}, \cmd{CONTACT}, \cmd{GIVE\_WAY}, \cmd{TURN}, \cmd{CLIMB}, \cmd{DESCEND}, \cmd{HEADING}, \cmd{SQUAWK},\\
&\cmd{HOLD\_POSITION}, \cmd{MONITOR}, \cmd{FOLLOW}, \cmd{TRAFFIC\_INFO}, \cmd{WAKE\_ADVISORY}, \cmd{ACK}, \cmd{REQUEST}\},
\end{align*}
즉 20종이다. 세 부분집합을 둔다.
\begin{itemize}
\item \textbf{복창 필수} $\Treq^{\rho}$ ($\KR$에 둔다): \textbf{지역 규정에 따라 정의한다}. US는 FAA JO 7110.65 §2-4-3 \cite{faa7110}을 근거로 등급을 매겼다. 본문 b의 관제사 의무(\texttt{required\_by\_atc}: \cmd{HOLD\_SHORT}), NOTE 1의 권고(\texttt{recommended\_runway}: 택시·이륙·착륙 허가의 활주로, \texttt{recommended\_numbers}: 주파수·고도·기수·코드), 규정에는 없지만 활주로 진입·정지 지시라 팀이 정한 \texttt{team\_rule}(\cmd{CROSS\_RUNWAY}, \cmd{LINE\_UP\_AND\_WAIT}, \cmd{HOLD\_POSITION})을 합쳐 13종이다. 응답만 하면 되는 \texttt{ack\_only}는 제외한다. KR/JP는 ICAO Doc 4444(PANS-ATM)의 복창 요건 \cite{icao4444}과 각국 AIP를 따르며, 현재 목록(12종, \cmd{MONITOR} 제외)은 원문 미검증이다 (RQ4 단계 과제).
\item \textbf{상태형}:
\begin{align*}\Tst = \{&\cmd{CLEARED\_TO\_LAND},\cmd{CLEARED\_FOR\_TAKEOFF},\cmd{LINE\_UP\_AND\_WAIT},\cmd{HOLD\_SHORT},\\ &\cmd{HOLD\_POSITION},\cmd{CROSS\_RUNWAY},\cmd{TAXI\_VIA},\cmd{CONTACT},\cmd{MONITOR}\}\end{align*}
\item \textbf{정보형} $\Tadv = \{\cmd{TRAFFIC\_INFO},\cmd{WAKE\_ADVISORY}\}$, $\Tadv \cap \Treq = \emptyset$ (\ch{C7}).
\end{itemize}

\textbf{[방법·근거]} 명령 구조와 이름은 DLR 공식 사양 v3.00 \cite{dlr2025}을 기준으로 하고, 미국 지상·타워 교신에 필요한 20종을 골라 정형화했다. 공개된 기계 판독 파일(OWL/JSON)이 없으므로, 사람이 정한 대응을 생성 코드가 PDF 원문과 자동 대조한다(조건부 허가는 범위 제외). 유럽 온톨로지로 미국 지상관제 교신을 표현했을 때 명령 종류 손실이 2.6\%(값 포함 4.6\%)라는 실증 \cite{helmke2025loss}이 이 선택을 뒷받침한다. 복창 규정은 FAA JO 7110.65 \cite{faa7110}, 미국 어휘는 MITRE--DLR 비교 \cite{chen2023}를 따른다. 노드 라벨은 atmonto IRI(예: \texttt{atm:Runway})에 매핑한다 \cite{keller2017,keller2019}. $\Sigma$는 SHACL shape로 기술한다 \cite{shacl}. 예: ``\cmd{Instruction}은 정확히 하나의 \cmd{Transmission}에 포함된다'', ``\cmd{CLEARED\_TO\_LAND}의 대상은 \cmd{Runway}이다''.

%----------------------------------------------------------
\paragraph{M1. STT (M1a + M1b).} 교신 음성을 송신 단위 전사 텍스트로 바꾼다. Whisper 경로는 M1a(음성 구간 분할) $\to$ M1b(전사)로, MOSS-ATC 경로는 한 모델이 분할·전사·화자 구분을 함께 한다. 어느 경로든 파일 이름의 녹음 맥락(공항·주파수·시작 시각)은 M1a가 풀어 송신마다 꼬리표로 붙인다.

\subsection{M1a. 송신 분할 (Transmission Segmentation)}
\textbf{[해결 문제]} \ch{C1} 무음 구간 환각의 원천 차단, \ch{C9} 클립 단위 맥락 부여.

\textbf{[입력]} 피드별 음성과 맥락 $\{(a^{\varphi},\Ctx^{\varphi})\}_{\varphi\in\Phi}$. $\Ctx^{\varphi}$는 파일명 \texttt{\{ICAO\}-\{Facility\}-\{Date\}-\{Time\}Z}에서 파싱한다. 시설 필터 $\varphi\in\Phi_{\mathrm{use}}$를 적용한다(\S\ref{sec:data}).

\textbf{[출력]} $S = \{s_i = [t^s_i, t^e_i)\}_{i=1}^{n}$, $x_i$, 각 $s_i$에 $\Ctx_i=\Ctx^{\varphi(i)}$를 부여. 모든 피드의 송신을 \textbf{절대 UTC 시각} $t = \tau_0^{\varphi} + n/f_s$ 기준으로 하나의 시간순 스트림으로 병합한다: $S=\mathrm{sort}_t\big(\bigcup_{\varphi} S^{\varphi}\big)$.

\textbf{[정의]} 프레임 $k$(25\,ms 창, 10\,ms 이동)의 에너지는
\begin{equation}
E_k = 10\log_{10}\Big(\frac{1}{L}\sum_{n\in\mathrm{frame}_k} a[n]^2\Big)\ [\mathrm{dB}]
\end{equation}
이고, 활성 프레임은
\begin{equation}
\mathrm{act}_k = \mathbb{1}\big[E_k > \theta_{\mathrm{dig}}\big]\cdot \mathbb{1}\big[p^{\mathrm{vad}}_k > \theta_{\mathrm{vad}}\big],\qquad \theta_{\mathrm{dig}}=-90\,\mathrm{dB}
\end{equation}
이다. $s_i$는 활성 프레임의 최대 연속 구간이다. 간격이 $< g_{\min}$이면 병합하고, 길이가 $< d_{\min}$이면 제거한다.

\textbf{[방법·근거]} 편집된 유튜브 샘플은 57.5\%가 디지털 무음이라 에너지 문턱이 효과적이었다. 반면 LiveATC 원본(KATL \cmd{Gnd-All})은 디지털 무음이 \textbf{0\%}로, 송신 사이가 잡음으로 채워져 있다. 따라서 원본에서는 \textbf{신경망 VAD(Silero VAD)를 주 판정}으로 하고, 에너지 문턱은 편집 음원에서만 보조로 쓴다. 무음 구간을 ASR에 넣지 않는 것이 환각을 막는 가장 직접적인 방법이다 \cite{koenecke2024}. 샘플에서 VAD 사용 시 환각이 사라짐을 확인했다.

%----------------------------------------------------------
\subsection{M1b. 맥락 바이어싱 ASR + 환각 필터}
\textbf{[해결 문제]} \ch{C2} 도메인 어휘 오인식, \ch{C1} 잔여 환각, \ch{C11} 라벨 없이 적응.

\textbf{[입력]} $x_i$, $\Vdom(\alpha_{\mathrm{apt}})$, 활성 콜사인 목록 $\mathcal{F}_{\mathrm{act}}(t_i) \subseteq V_{\cmd{Flight}}(\G_{<t_i})$.

\textbf{[출력]} $(w_i, c_i, \mathrm{lang}_i)$. $w_i\in\Sigma^*$이고, $c_i = (\bar\ell_i,\ p^{\mathrm{ns}}_i)$는 평균 토큰 로그확률과 무음 확률, $\mathrm{lang}_i$는 Whisper 언어 식별 결과다(RQ4). 선택적으로 $N$-best $\{(w_i^{(k)},P^{(k)})\}_{k=1}^{K}$를 낸다.

\textbf{[정의]}
\begin{align}
w_i &= \arg\max_{w}\; P_\theta\big(w \mid x_i,\ \mathrm{prompt}(\Vdom,\mathcal{F}_{\mathrm{act}}(t_i))\big),\\
\mathrm{keep}(i) &= \mathbb{1}\big[p^{\mathrm{ns}}_i < \delta_{\mathrm{ns}}\big]\cdot\mathbb{1}\big[\bar\ell_i > \delta_{\ell}\big]\cdot\mathbb{1}\big[w_i\notin\mathcal{B}\big],
\end{align}
여기서 $\mathcal{B}$는 알려진 환각 문구 집합이다 (예: \texttt{"thanks for watching"}, \texttt{"you"}).
\emph{비영어 발화 처리 (\ch{C14})}: $\mathrm{lang}_i\neq\texttt{en}$이고 언어 확률이 $\delta_{\mathrm{lang}}$ 이상이면, 송신 노드를 \cmd{UNSUPPORTED\_LANG}으로 표시해 \textbf{그래프에는 남기되} M2a--M3 처리는 건너뛴다. 출처 보존(P2)을 지키면서, 자국어 혼용 빈도를 RQ4 실패 원인으로 집계할 수 있다.

\textbf{[방법·근거]} Whisper large-v3 \cite{whisper}에 초기 프롬프트로 도메인 어휘와 활성 콜사인을 넣는다. 감시 데이터로 콜사인 n-gram을 부스팅해 콜사인 인식을 크게 개선한 연구 \cite{nigmatulina2022}의 \emph{라벨 없이 쓸 수 있는} 경량 근사다. 그래프 $\G_{<t_i}$에서 활성 콜사인을 공급하므로 KG가 ASR을 돕는 되먹임이 생긴다. ATC 미세조정 모델 \cite{vandoorn}은 유럽 데이터로 학습되어 미국 억양에서 성능이 떨어진다는 보고가 있어 \cite{usatc2025}, 비교 실험으로만 둔다.

%----------------------------------------------------------
\paragraph{M2. 지시 추출 (M2a + M2b).} 송신 텍스트에서 화자 역할(관제사/조종사)을 정하고 지시를 온톨로지 형식으로 뽑는다. M2a가 정한 화자 단위 역할이 최종 역할이다. 화자 번호가 없는 Whisper 경로에서만 M2b의 LLM이 규칙 판정을 보완한다(규칙만으로 역할 미정 51\% $\to$ 0.2\%).

\subsection{M2a. 화자 역할 판정}
\textbf{[해결 문제]} \ch{C3} 화자 역할 불명. 송신 텍스트만 보면 ``Delta 123''으로 시작하는 말이 관제사인지 조종사인지 정하기 어려운 경우가 많다(규칙만으로는 51\% 미정).

\textbf{[입력]} 한 녹음 파일의 송신 목록 $\{(w_i,\sigma_i)\}_{i}$. $\sigma_i$는 M1(MOSS)이 붙인 화자 번호(S01, S02 \dots)로, 같은 파일 안에서 같은 목소리면 같은 번호다. 온톨로지는 넣지 않는다(역할 판정에 명령 어휘 목록이 필요 없다).

\textbf{[출력]} 화자 역할 표 $\rho:\sigma\mapsto r\in\{\ATCO,\PILOT,\texttt{UNKNOWN}\}$ (\texttt{speaker\_roles.json}: 화자별 역할·송신 수·근거)와 송신별 역할 $r_i=\rho(\sigma_i)$ (\texttt{turns.jsonl}: \texttt{tx\_id}, \texttt{speaker}, \texttt{role}, \texttt{role\_source}, 텍스트·시각). \texttt{role\_source}는 \texttt{llm\_speaker}(LLM 판정), \texttt{rule\_override}(LLM 판정을 콜사인 위치로 교정), \texttt{rule\_fallback}(LLM 답이 없어 규칙 다수결) 중 하나다.

\textbf{[정의]}
\begin{equation}
\rho(\sigma) = \arg\max_{r}\; P_{\mathrm{LLM}}\big(r \mid W_\sigma,\ |\mathcal{C}_\sigma|\big),\qquad
W_\sigma=\{w_i:\sigma_i=\sigma\},\quad \mathcal{C}_\sigma=\textstyle\bigcup_{i:\sigma_i=\sigma}\mathrm{callsigns}(w_i)
\end{equation}
LLM에는 화자마다 송신 수, 부른 콜사인 종류 수 $|\mathcal{C}_\sigma|$, 긴 송신 5개를 준다. 판정 근거는 FAA JO 7110.65BB의 통신 형식이다. 관제사는 여러 항공기를 부르며 콜사인을 앞에 두고 지시한다(\S2-4-8: 항공기 식별부호를 먼저). 조종사는 자기 콜사인 하나만 쓰고 복창 끝에 붙인다. 실제 예(\texttt{KATL-Gnd-All-Aug-26-2024-0000Z}): S06은 송신 157건에서 콜사인 60종을 부르고 콜사인이 앞에 온 송신이 82건, 뒤에 온 송신이 2건이라 \ATCO. S81은 13건 중 10건이 ``nine left Mike two \dots\ Delta twenty four twenty five''처럼 콜사인을 끝에 붙인 복창이라 \PILOT.

\textbf{[대체 규칙]} LLM 답에 화자가 빠지거나 형식이 틀리면 그 화자의 송신마다 규칙 특징 $\phi(w)=(\mathrm{pos}_{\mathrm{cs}}(w),\ \mathbb{1}[\text{시설 호출로 시작}],\ \mathbb{1}[\text{지시 동사}],\dots)$로 판정해 다수결한다. 콜사인을 3종 이상 부른 화자는 \ATCO 쪽으로 판정한다. 또 LLM이 \ATCO{}라고 답했어도 콜사인이 뒤에 온 송신이 앞에 온 송신보다 많으면(2건 이상) \PILOT{}으로 바꾼다(\texttt{rule\_override}). Qwen3-4B는 복창에도 지시 어휘(taxi, hold short)가 들어 있어 조종사를 관제사로 고르는 오류가 잦았다.

\textbf{[방법·근거]} 송신 하나보다 화자 하나에 모인 송신 여러 개가 역할을 더 잘 드러낸다. 텍스트 기반 역할 판정이 ATC에서 효과적이라는 결과 \cite{bertraffic}를 따르되, 판정 단위를 화자로 올렸다. MOSS가 화자를 구분하므로 한 송신 안에서 화자가 바뀌는 경우(턴 분리)는 다루지 않는다. 한계: MOSS는 목소리가 비슷한 조종사 여러 명을 한 번호로 묶기도 한다(위 S81은 서로 다른 콜사인 6종). 이때 한 번호 안의 송신은 모두 같은 역할이라 판정에는 문제가 없지만, 관제사와 조종사가 한 번호로 묶이면 그 화자의 일부 송신은 역할이 틀린다. 또 MOSS는 숫자를 단어로 적으므로(``twenty four twenty five'') 콜사인을 찾기 전에 묶음 숫자를 숫자로 바꾼다(\texttt{common/atc\_text.py: group\_words\_to\_digits}). 화자 번호가 없는 Whisper 경로에서는 기존처럼 송신마다 규칙으로 판정하고, 역할이 바뀌는 문장 경계에서 자른다(KATL에서 20,022건 중 61건).

%----------------------------------------------------------
\subsection{M2b. 온톨로지 기반 LLM 의미 추출}
\textbf{[해결 문제]} \ch{C6} 개념 단위 표현, \ch{C7} 지시와 정보 구분, \ch{C11} 라벨 부족, \ch{C13} 형식 오류.

\textbf{[입력]} $w_i$, $\Ont$(\texttt{command\_types}, \texttt{value\_formats}), few-shot 예시 $\mathcal{E}=\{(w^{(e)},u^{(e)})\}_{e=1}^{k}$.

\textbf{[출력]}
\begin{equation}
u_i = \big(\hat c_i,\; I_i\big),\qquad I_i=\{\iota_{ij}\}_{j},\quad \iota=(\kappa, v, \mu, q, g),\ \kappa\in\Tcmd,
\end{equation}
\begin{itemize}
\item $\hat c_i$: 콜사인 표면 문자열. 생략되면 $\varepsilon$.
\item $v$: 값, $\mu$: 단위.
\item $q$: 한정어 (예: \cmd{LEFT}, \cmd{AT}:\texttt{Juliet}).
\item $g$: 대상 개체 언급 (예: \texttt{"runway two two right"}).
\end{itemize}

\textbf{[정의]}
\begin{equation}
u_i = \arg\max_{u\,\in\,\mathcal{L}(\mathrm{JSchema}(\Ont))}\; P_{\mathrm{LLM}}\big(u \mid \mathrm{prompt}(w_i, r_i, \Ont, \mathcal{E})\big),
\end{equation}
여기서 $\mathcal{L}(\mathrm{JSchema}(\Ont))$는 $\Ont$에서 생성한 JSON Schema가 허용하는 문자열 집합이다. $\kappa$는 열거형(enum)으로 강제된다.

\textbf{[방법·근거]}
\begin{itemize}
\item \textbf{스키마 기반 추출(D2)}: 비교 가능한 정규 개념이 필요하다 \cite{feng2024,llmkgc2025}.
\item \textbf{LLM few-shot(D3)}: 대량 주석 없이 ATC 명령 추출에서 높은 성능이 보고되었다 \cite{mekerishvili2025}.
\item \textbf{제약 디코딩(D4)}: 형식 오류를 원천 차단한다 \cite{willard2023}.
\item \textbf{입력 근거 제약}: 대상 언급 $g$는 $w_i$의 부분 문자열이어야 한다는 제약을 프롬프트와 사후 검사로 둔다 (입력에 없는 개체 생성 방지, \ch{C13}) \cite{itext2kg}.
\item \textbf{LLM 백엔드 교체 가능 (\ch{C15})}: $P_{\mathrm{LLM}}\in\{P_{\mathrm{API}},\ P_{\mathrm{local}}\}$. 보안 자료(KR-2안)는 반드시 로컬 모델(RTX 4090에서 구동)로만 처리한다. 지역 간 비교가 LLM 차이로 오염되지 않도록, RQ4의 모든 조건은 \textbf{같은 백엔드}로 실행한다. 로컬 백엔드를 쓰면 US 검증 세트에도 같은 로컬 모델을 돌려 기준선을 맞춘다.
\end{itemize}

%----------------------------------------------------------
\subsection{M3. 정규화·Grounding·개체 해소}
\textbf{[해결 문제]} \ch{C4} 콜사인 오인식·생략, \ch{C5} 표기 변이, \ch{C9} 정적 지식 검증, \ch{C2} 잔여 오인식 교정.

\textbf{[입력]} $u_i$, $\KS$, $\G_{<t_i}$, $\Ctx_i$.

\textbf{[출력]}
\begin{itemize}
\item $\tilde u_i = (f_i,\ \tilde I_i,\ \mathcal{A}_i)$.
\item $f_i\in V_{\cmd{Flight}}\cup\{\bot\}$: 해석된 항공편.
\item $\tilde\iota = (\kappa,\ \nu(v),\ \mu,\ q,\ \ell(g))$, 여기서 $\ell(g)\in V_S\cup\{\bot\}$.
\item $\mathcal{A}_i$: 교정·연결 실패 기록 (예: \cmd{ASR\_CORRECTED}, \cmd{KB\_VIOLATION}).
\end{itemize}

\textbf{[정의]}
\begin{itemize}
\item \emph{값 정규화}: $\nu$는 숫자·음성 알파벳 정규화 함수다. 예: $\nu(\texttt{"two two right"}) = \nu(\texttt{"22 right"}) = \cmd{22R}$.
\item \emph{대상 연결}:
\begin{equation}
\ell(g) = \begin{cases}\arg\max_{v\in V_S,\ v\sqsubseteq \alpha_{\mathrm{apt}}}\ \mathrm{sim}(\nu(g), \mathrm{name}(v)) & \text{if } \max \ge \theta_\ell\\ \bot & \text{otherwise (}\cmd{KB\_VIOLATION}\text{)}\end{cases}
\end{equation}
\item \emph{콜사인 해소}: 활성 항공편 집합을 $\mathcal{F}_{\mathrm{act}}(t_i) = \{f : \exists x\in\G_{<t_i},\ \mathrm{flight}(x)=f,\ t_i - t(x) \le W\}$로 둔다.
\begin{equation}
f_i = \arg\max_{f\in\mathcal{F}_{\mathrm{act}}(t_i)\cup\{f_{\mathrm{new}}\}}\;\Big[\beta\,\mathrm{sim}_{\mathrm{ph}}\big(\hat c_i, \mathrm{tel}(f)\big) + (1-\beta)\,\mathrm{ctx}(f, i)\Big],
\end{equation}
  \begin{itemize}
  \item $\mathrm{sim}_{\mathrm{ph}}$: 발음 기반 유사도. 항공사 telephony와 편명 숫자를 따로 비교한다.
  \item $\mathrm{ctx}(f,i)$: 문맥 점수. $i$가 $f$에게 내려진 직전 지시의 개념을 반복하면(복창 후보) 높다.
  \item $f_{\mathrm{new}}$는 $\mathrm{tel}^{-1}(\hat c_i)$로 새로 만든 항공편이다. $\hat c_i=\varepsilon$이면 $\mathrm{ctx}$만으로 결정한다.
  \end{itemize}
\end{itemize}

\textbf{[방법·근거]}
\begin{itemize}
\item 맥락 정보(감시 데이터·활성 콜사인)를 이용한 콜사인 교정 \cite{nigmatulina2022,helmke2021}.
\item EDC의 사후 정규화(canonicalization) 단계 \cite{edc2024}.
\item iText2KG의 증분 개체 해소 \cite{itext2kg}.
\end{itemize}
예: \texttt{Alpha 2398}에서 \texttt{Alpha}는 telephony 사전에 없고 편명 \texttt{2398}이 일치하며, 직후 복창이 \texttt{Delta 2398}이다. 따라서 $f=\cmd{DAL2398}$로 해소하고 $\mathcal{A}_i \ni \cmd{ASR\_CORRECTED}$를 남긴다.

%----------------------------------------------------------
\subsection{M4. 증분 사건 중심 Temporal KG 구축}
\textbf{[해결 문제]} \ch{C8} 발화 간 관계, \ch{C10} 시간 유효성, \ch{C12} 출처·신뢰도, \ch{C13} 스키마 검증.

\textbf{[입력]} $\G_{t_{i-1}}$, $\tilde u_i$, $(s_i, r_i, w_i, c_i)$, $\Sigma$, $\KR$, $\Cfg$.

\textbf{[출력]} $\G_{t_i} = \G_{t_{i-1}} \oplus \Delta_i$.

\textbf{[정의] 사건 reification.} 송신 노드 $n^{T}_i$, 지시 노드 $n^{I}_{ij}$를 만들고 다음 간선을 추가한다.
\begin{align}
\Delta_i = \{\,& (n^{T}_i,\cmd{SPOKEN\_BY},r_i),\ (n^{T}_i,\cmd{ADDRESSES},f_i),\ (n^{T}_i,\cmd{CONTAINS},n^{I}_{ij}),\nonumber\\
& (n^{I}_{ij},\cmd{TARGET},\ell(g_{ij})),\ (n^{I}_{ij},\cmd{NEXT},n^{I}_{\mathrm{prev}(f_i)})\,\} \;\cup\; \Delta^{\mathrm{rb}}_i \;\cup\; \Delta^{\mathrm{inv}}_i .
\end{align}

\textbf{[정의] 이중 시간 (bi-temporal).} 모든 동적 원소 $x$에 대해 다음을 둔다.
\begin{itemize}
\item 사건 시간 $\tau(x) = [t^{\mathrm{from}}_x, t^{\mathrm{to}}_x)$: 사실이 실제로 유효한 구간.
\item 처리 시간 $t'_x$: 시스템에 적재된 시각.
\item $\mathrm{prov}(x)=\{s_i\}$, $\mathrm{conf}(x) = g(c_i, \mathrm{conf}_{\mathrm{LLM}}, \mathcal{A}_i)$.
\end{itemize}
상태형 지시 $\iota\in\Tst$의 유효 구간은
\begin{equation}
\tau(\iota) = \big[t^s_i,\ t^{\mathrm{to}}\big),\quad
t^{\mathrm{to}} = \min\Big(\{t^s_k : k>i,\ f_k=f_i,\ \exists\iota'\in \tilde I_k,\ \mathrm{sup}(\iota',\iota)\} \cup \{t^s_i + T_{\max}(\kappa)\}\Big)
\end{equation}
이다. $\mathrm{sup}(\iota',\iota)$는 대체(supersede) 관계로, $\KR$(\texttt{rules.json})에 표로 정의한다. 예: \cmd{CROSS\_RUNWAY}는 \cmd{HOLD\_SHORT}·\cmd{HOLD\_POSITION}을, \cmd{CLEARED\_FOR\_TAKEOFF}는 \cmd{LINE\_UP\_AND\_WAIT}·\cmd{HOLD\_SHORT}를 종료하고, \cmd{CONTACT}·\cmd{MONITOR}는 그 항공편의 모든 상태를 종료한다. 이 관계는 규정 문장이 아니라 지시의 의미에서 정했다. $T_{\max}(\kappa)$는 대체 지시가 녹음에 없을 때의 상한으로 $\Cfg$(\texttt{kg\_build.json})에 둔다(예: \cmd{HOLD\_SHORT} 900\,s, \cmd{CROSS\_RUNWAY} 180\,s, 기본 600\,s). KATL 전체 KG에서 관제사 상태형 지시 2{,}993개 중 2{,}251개가 이 상한으로 종료되었다. 비상태형 지시는 $\tau(\iota)=[t^s_i,t^e_i)$이다.

\textbf{삭제하지 않고 무효화한다}: 대체가 일어나면 $\Delta^{\mathrm{inv}}_i$가 기존 간선의 $t^{\mathrm{to}}$만 설정하고 간선은 남긴다. 따라서 과거 시점 질의가 가능하다 \cite{zep2025}.

\textbf{[정의] 복창 연결.}
\begin{equation}
\Delta^{\mathrm{rb}}_i = \Big\{(n^T_i,\cmd{READBACK\_OF},n^T_j)\ :\ r_i=\PILOT,\ r_j=\ATCO,\ f_i=f_j,\ 0 < t^s_i - t^e_j \le \Delta_{\mathrm{rb}},\ j=\max\{j'<i\}\Big\},
\end{equation}
$f_i=\bot$이면 직전 관제사 송신의 수신자와의 개념 중복도 $\mathrm{ovl}(\tilde I_i,\tilde I_j)\ge\theta_{\mathrm{ovl}}$로 대신 판정한다. $\Delta_{\mathrm{rb}}=20$\,s는 $\Cfg$에 둔다. KATL에서 연결된 복창 404건 중 351건이 5\,s 이내, 95\%가 12.4\,s 이내였다(20\,s 창 안에서 측정).

\textbf{[정의] 검증.} 적재 후 $\G_{t_i}\models\Sigma$를 검사한다. 위반하면 해당 $\Delta$를 \cmd{quarantine} 라벨로 격리한다.

\textbf{[방법·근거]}
\begin{itemize}
\item n-항 관계 reification \cite{naryw3c}과 사건 중심 TKG \cite{knez2023}.
\item 이중 시간·무효화 모델 \cite{zep2025}.
\item 관측 시간과 유효 시간의 분리 \cite{atom2025}.
\item 시간 관계는 Allen 구간 대수 \cite{allen1983}와 OWL-Time \cite{owltime} 의미로 해석한다.
\end{itemize}

%----------------------------------------------------------
\subsection{M5. 규칙 기반 추론과 이상 탐지}
\textbf{[해결 문제]} \ch{C6} 개념 단위 복창 검증, \ch{C4} 콜사인 누락, \ch{C7} 정보형 제외, \ch{C12} 근거 제시.

\textbf{[입력]} $\G_{t}$, 규칙 집합 $\mathcal{R}=\{\rho_1,\dots,\rho_7\}$, $\KR$의 $\Treq^{\rho}$와 $\mathcal{T}_{\mathrm{rwy}}$.

\textbf{[출력]} $\Yset=\{(y_k,\pi_k,\gamma_k)\}$와 \cmd{READBACK\_OF} 간선의 상태 속성.

\textbf{[정의] 개념 단위 복창 비교.} 지시 송신 $j$와 복창 송신 $i$에 대해 다음을 정의한다.
\begin{align}
\mathrm{Req}(j) &= \{\tilde\iota\in\tilde I_j : \kappa(\tilde\iota)\in\Treq\},\\
m(\tilde\iota,\tilde\iota') &= \mathbb{1}\big[\kappa=\kappa'\ \wedge\ \ell(g)=\ell(g')\ \wedge\ \nu(v)=\nu(v')\ \wedge\ q\simeq q'\big],\\
\mathrm{Miss}(j,i) &= \{\tilde\iota\in\mathrm{Req}(j) : \nexists\,\tilde\iota'\in\tilde I_i,\ \kappa(\tilde\iota')=\kappa(\tilde\iota)\},\\
\mathrm{Conf}(j,i) &= \{\tilde\iota\in\mathrm{Req}(j) : \exists\,\tilde\iota'\in\tilde I_i,\ \kappa'=\kappa,\ m(\tilde\iota,\tilde\iota')=0\}.
\end{align}
상태 판정은
\begin{equation}
\mathrm{status}(j,i)=\begin{cases}
\cmd{MISMATCH} & \mathrm{Conf}(j,i)\neq\emptyset\\
\cmd{PARTIAL} & \mathrm{Conf}=\emptyset,\ \mathrm{Miss}\neq\emptyset\\
\cmd{MATCH} & \text{otherwise}
\end{cases}
\end{equation}
이다. $\Tadv$는 $\Treq$에 속하지 않으므로 자동으로 비교에서 빠진다 (\ch{C7}).

\textbf{규칙 목록}
\begin{itemize}
\item[$\rho_1$] \cmd{MISMATCH}/\cmd{PARTIAL}: 위 판정. 예: 좌/우 활주로 혼동은 $\ell(g)\neq\ell(g')$로 잡힌다 \cite{cardosi94}.
\item[$\rho_2$] \cmd{NO\_READBACK}: $\mathrm{Req}(j)\neq\emptyset$이고 $\nexists i$ with $(i,\cmd{READBACK\_OF},j)$.
\item[$\rho_3$] \cmd{NO\_CALLSIGN}: $(i,\cmd{READBACK\_OF},j)\wedge \hat c_i=\varepsilon$ (FAA JO 7110.65 §2-4-3 \cite{faa7110}).
\item[$\rho_4$] \cmd{RWY\_OVERLAP}: 다음을 만족하는 $\tilde\iota_a,\tilde\iota_b$가 존재 (Allen 관계 \cite{allen1983}).
\[\kappa_a,\kappa_b\in\mathcal{T}_{\mathrm{rwy}},\quad \ell(g_a)=\ell(g_b),\quad f_a\neq f_b,\quad \tau(\tilde\iota_a)\ \mathsf{overlaps}\ \tau(\tilde\iota_b),\]
\[\mathcal{T}_{\mathrm{rwy}}=\{\cmd{CLEARED\_TO\_LAND},\cmd{CLEARED\_FOR\_TAKEOFF},\cmd{LINE\_UP\_AND\_WAIT},\cmd{CROSS\_RUNWAY}\}.\] 실제 운영에서는 예상 분리 허가 등이 허용되므로 위반이 아니라 \emph{주의 힌트}로 표시한다.
\item[$\rho_5$] \cmd{POST\_HANDOFF}: $(f,\cmd{CONTACT})$의 $t^{\mathrm{from}}$ 이후 같은 위치 $\phi_{\mathrm{pos}}$에서 $f$에게 새 지시가 나감.
\item[$\rho_6$] \cmd{KB\_VIOLATION}: $\ell(g)=\bot$. 공항에 없는 활주로·유도로이며, ASR 오류 또는 맥락 오설정을 시사한다.
\end{itemize}

\textbf{[근거 경로와 신뢰도]} 각 이상에 근거 부분그래프 $\pi_k$를 붙인다. $\pi_k$는 관련 \cmd{Transmission}, \cmd{Instruction}, 대상 노드와 그 간선으로 구성된다. 신뢰도는
\begin{equation}
\gamma_k = \min_{x\in\pi_k}\mathrm{conf}(x)\cdot\big(1-\lambda\,\mathbb{1}[\cmd{ASR\_CORRECTED}\in\mathcal{A}]\big)
\end{equation}
이다. ASR 오류로 인한 거짓 경보와 실제 복창 오류를 사용자가 구분할 수 있게 한다 (\ch{C12}).

\textbf{[방법·근거]} 개념 수준 비교는 복창 오류 탐지 연구의 핵심 결론이다 \cite{helmke2021}. 규칙은 규정에서 직접 도출하므로 학습 데이터가 필요 없고 설명 가능하다 (D8). 기호적 KG와 LLM의 상호 보완 관점 \cite{pan2024}에 따라 LLM은 추출(\md{M2b})에만 쓰고 판정은 기호 규칙으로 한다.

%----------------------------------------------------------
\subsection{M6. 시점 질의와 시각화}
\textbf{[해결 문제]} \ch{C10} 시점 기반 상황 인식, \ch{C12} 근거 시각화.

\textbf{[입력]} $\G$, 시각 $t$, 질의 $q\in\mathcal{Q}$.

\textbf{[출력]} $Q(\G,t,q)$. 답 집합과 근거 부분그래프.

\textbf{[정의]} 시점 스냅샷과 질의 템플릿은 다음과 같다.
\begin{equation}
\G|_t = \big(V,\ \{e\in E : t\in\tau(e)\}\big),\qquad Q(\G,t,q) = \mathrm{eval}_q(\G|_t).
\end{equation}
\begin{itemize}
\item $q_1$: $\{f : \exists\iota,\ (f,\iota)\in\G|_t,\ \kappa(\iota)=\cmd{CLEARED\_TO\_LAND},\ \ell(g)=\cmd{28L}\}$ (지금 28L 착륙 허가 항공편. KATL처럼 Tower 피드가 없으면 $\kappa=\cmd{CROSS\_RUNWAY}$/\cmd{HOLD\_SHORT}로 바꿔 ``지금 8R 횡단 허가·대기 중인 항공편''으로 대체)
\item $q_2$: $\mathrm{timeline}(f) = \langle \iota_1 \prec \iota_2 \prec \dots\rangle$ (\cmd{NEXT} 경로)
\item $q_3$: $\{(y,\pi)\in\Yset : t-\Delta \le t(y) \le t\}$ (최근 이상과 근거)
\end{itemize}

\textbf{[방법·근거]} Cypher 시간 필터로 구현하고 Neo4j Browser로 시각화한다. Level 3 ``예측''은 활성 상태 구간의 잔여 시간 정도로 제한한다. TKG 기반 미래 예측 \cite{cai2024tkg}은 범위 밖이다 (D7).

%----------------------------------------------------------
\subsection{M7. 파이프라인 진단 하네스 (RQ4 전용)}
\textbf{[해결 문제]} \ch{C14} 지역 전이의 원인 분석, \ch{C15} 보고 범위 통제, \ch{C12} 오류 전파 추적.

\textbf{[입력]} 모든 모듈의 중간 산출물 로그 $\mathcal{J}=\{(i, k, \mathrm{out}_{\mathrm{M}k}(i))\}_{k=1}^{7}$, 소량 정답 $\{\mathrm{gold}_{\mathrm{M}k}(i)\}$. 정답은 국가별 15--20개 송신에만 단다.

\textbf{[출력]} 지역·조건별 지표 $(\mathrm{CSR}_\rho,\ \mathrm{CmdF1}_\rho)$와 처음 실패한 모듈의 분포 $P_\rho(m^\dagger)$ (\S\ref{sec:rq4}). 보안 자료에 대해서는 \textbf{집계 수치만} 내보낸다.

\textbf{[정의]} 각 모듈 출력을 송신 ID $i$를 키로 기록하므로, $\mathrm{prov}$(P2)와 같은 키로 추적할 수 있다. 정답 비교 함수는 M5의 개념 일치 $m(\cdot,\cdot)$을 그대로 재사용한다.

\textbf{[방법·근거]} KG 자체의 정량 평가가 아니라 \emph{모듈 단위 오류 귀속}이다. 범위 결정(KG 평가 생략)과 충돌하지 않으면서, 지역 간 비교의 원인을 설명한다. 계단식 구조(D1)이기 때문에 이런 귀속이 가능하다.

\clearpage
%==========================================================
\section{KG 스키마와 시간 모델}
%==========================================================
\begin{figure}[h!]
\centering
\resizebox{0.95\textwidth}{!}{%
\begin{tikzpicture}[
  font=\small,
  dyn/.style={draw, thick, rounded corners=2pt, fill=comp, minimum height=0.8cm, minimum width=2.4cm, align=center},
  sta/.style={draw, thick, fill=res, minimum height=0.8cm, minimum width=2.2cm, align=center},
  e/.style={-{Stealth[length=2.2mm]}, thick},
  el/.style={font=\scriptsize\sffamily, fill=white, inner sep=1pt}
]
\node[dyn] (tx) at (0,0) {\textbf{Transmission}\\{\scriptsize $s_i, w_i, c_i, t'$}};
\node[dyn] (tx2) at (0,-3.2) {\textbf{Transmission}\\{\scriptsize (PILOT)}};
\node[sta] (role) at (-4.2,0) {\textbf{Role}\\{\scriptsize ATCO / PILOT}};
\node[dyn] (fl) at (4.6,0) {\textbf{Flight}\\{\scriptsize callsign}};
\node[sta] (al) at (8.6,0) {\textbf{Airline}\\{\scriptsize icao, telephony}};
\node[dyn] (ins) at (4.6,-3.2) {\textbf{Instruction}\\{\scriptsize $\kappa,v,\mu,q,\tau,\mathrm{conf}$}};
\node[sta] (rw) at (8.6,-3.2) {\textbf{Runway / Taxiway}\\{\scriptsize / Facility}};
\node[sta] (ap) at (8.6,-5.6) {\textbf{Airport}};
\node[dyn] (cm) at (4.6,2.2) {\textbf{CallsignMention}\\{\scriptsize $\hat c_i$ (원 표면형)}};
\node[dyn] (ins2) at (1.4,-5.6) {\textbf{Instruction}\\{\scriptsize (다음)}};

\draw[e] (tx) -- node[el,above]{SPOKEN\_BY} (role);
\draw[e] (tx) -- node[el,above]{ADDRESSES} (fl);
\draw[e] (fl) -- node[el,above]{OPERATED\_BY} (al);
\draw[e] (tx) -- node[el,sloped,above]{CONTAINS} (ins);
\draw[e] (ins) -- node[el,above]{TARGET} (rw);
\draw[e] (rw) -- node[el,right]{PART\_OF} (ap);
\draw[e] (fl) -- node[el,right]{ALIAS} (cm);
\draw[e, red!60!black] (tx2) -- node[el,left,text=red!60!black]{READBACK\_OF \{status\}} (tx);
\draw[e] (ins) -- node[el,sloped,below]{NEXT \{$\Delta t$\}} (ins2);
\draw[e, dashed] (tx2) -- node[el,below]{CONTAINS} (ins2);

\node[font=\scriptsize, align=left, anchor=north west] at (-5.4,-4.4) {
\tikz\node[fill=comp,draw,minimum width=4mm,minimum height=3mm]{}; 동적(사건) 계층: $\tau,\mathrm{prov},\mathrm{conf}$ 보유\\[2pt]
\tikz\node[fill=res,draw,minimum width=4mm,minimum height=3mm]{}; 정적 계층 $\KS$ (atmonto IRI 매핑)};
\end{tikzpicture}}
\caption{KG 스키마. 송신과 지시를 사건 노드로 reify하고(\ch{C8}), 오인식 표면형은 \cmd{CallsignMention}으로 보존한다(\ch{C4}, \ch{C12}).}
\label{fig:schema}
\end{figure}

\begin{figure}[h!]
\centering
\resizebox{\textwidth}{!}{%
\begin{tikzpicture}[font=\small, x=0.072cm, y=1cm]
% axis: sample seconds 170..250 -> scaled
\draw[-{Stealth}, thick] (170,0) -- (252,0) node[right]{$t$ [s]};
\foreach \x in {170,180,...,250} { \draw (\x,0.08) -- (\x,-0.08) node[below,font=\scriptsize]{\x}; }
% transmissions (sample, SFO part)
\foreach \a/\b/\lab/\yy in {177.7/182.1/{ATCO: AAL2069 CTL 28R}/1.0, 182.5/185.6/{PILOT: rb AAL2069}/1.55, 188.9/193.1/{PILOT: UAL1947? 28L}/2.1, 193.2/199.9/{ATCO: CTL 28L}/2.65, 200.4/205.1/{ATCO: ANZ7 CONTACT NorCal}/3.2, 241.0/246.3/{ATCO: AAL276 TAXI A A1}/1.0, 246.5/247.8/{PILOT: rb (no cs)}/1.55} {
  \fill[blue!25] (\a,\yy-0.17) rectangle (\b,\yy+0.17);
  \node[anchor=west,font=\scriptsize] at (\b+0.6,\yy) {\lab};
}
% validity bars
\fill[green!45!black!60] (177.7,-0.9) rectangle (250,-0.65);
\node[anchor=east,font=\scriptsize] at (176,-0.78) {$\tau$(CTL 28R, AAL2069)};
\fill[orange!70] (193.2,-1.45) rectangle (250,-1.2);
\node[anchor=east,font=\scriptsize] at (191,-1.33) {$\tau$(CTL 28L, UAL1947?)};
\draw[red!60!black, thick, -{Stealth}] (184,1.35) -- (181,1.18);
\node[font=\scriptsize, red!60!black, anchor=east] at (176,1.55) {READBACK\_OF: MATCH};
\draw[red!60!black, thick, -{Stealth}] (252,1.85) -- (247.2,1.75);
\node[font=\scriptsize, red!60!black, anchor=west] at (249,2.0) {$\rho_3$: NO\_CALLSIGN?};
\draw[dashed] (230,-1.6) -- (230,0.45) node[above,font=\scriptsize, fill=white, inner sep=1pt]{$\G|_{230}$};
\node[font=\scriptsize, align=left, anchor=west] at (231,-2.1) {$q_1(\G|_{230})$: 28R$\to$AAL2069, 28L$\to$UAL1947?\\ $\ell(g)$ 상이 $\Rightarrow$ $\rho_4$ 미발동};
\end{tikzpicture}}
\caption{샘플(SFO 구간)의 시간 모델. 위는 송신 구간, 아래는 상태형 지시의 유효 구간 $\tau$다. 유효 구간 덕분에 시점 스냅샷 $\G|_t$에서 ``지금'' 유효한 허가만 질의된다(\ch{C10}). 화자·콜사인 표기는 전사 기반 추정치다.}
\label{fig:time}
\end{figure}

%==========================================================
\section{추적성 행렬: 문제정의 $\leftrightarrow$ 모듈}\label{sec:trace}
%==========================================================
$\bullet$ = 주 해결 모듈, $\circ$ = 보조 모듈. 모든 행에 $\bullet$이 최소 1개 있으므로 정의한 문제가 빠짐없이 다뤄진다. 모든 열에도 $\bullet$이 최소 1개 있으므로 근거 없는 모듈이 없다.

\begin{center}\small
\setlength{\tabcolsep}{5pt}
\begin{tabular}{@{}l l c c c c c c c c c c@{}}
\toprule
& 문제 & M0 & M1a & M1b & M2a & M2b & M3 & M4 & M5 & M6\& M7 \\midrule
\multirow{4}{*}{\rotatebox{90}{\scriptsize RQ}}
& RQ1 표현 (지시–복창–이관 + 유효시간) & $\bullet$ &  &  & $\circ$ & $\bullet$ &  & $\bullet$ &  & &  \\
& RQ2 교정·연결 (STT 오류 탐지·교정) & $\circ$ &  & $\bullet$ &  &  & $\bullet$ & $\circ$ & $\circ$ & &  \\
& RQ3 활용 (복창 오류·상황 인식) &  &  &  &  &  &  & $\circ$ & $\bullet$ & $\bullet$&  \\
& RQ4 일반화 (한국·일본 교신 적용) & $\bullet$ &  & $\circ$ &  & $\circ$ & $\circ$ &  & $\circ$ & \& $\bullet$ \\midrule
& \ch{C1} 무음 환각 &  & $\bullet$ & $\bullet$ &  &  &  &  &  & &  \\
& \ch{C2} 협대역·도메인 어휘 & $\circ$ &  & $\bullet$ &  &  & $\bullet$ &  &  & &  \\
& \ch{C3} 화자 경계·역할 &  & $\circ$ &  & $\bullet$ &  &  &  &  & &  \\
& \ch{C4} 콜사인 오인식·생략 &  &  & $\circ$ &  &  & $\bullet$ &  & $\bullet$ & &  \\
& \ch{C5} 표기 변이 &  &  &  &  &  & $\bullet$ &  &  & &  \\
& \ch{C6} 바꿔 말하는 복창 & $\circ$ &  &  &  & $\bullet$ &  &  & $\bullet$ & &  \\
& \ch{C7} 지시/정보 구분 & $\bullet$ &  &  &  & $\bullet$ &  &  & $\circ$ & &  \\
& \ch{C8} 발화 간 관계 &  &  &  &  &  &  & $\bullet$ & $\circ$ & &  \\
& \ch{C9} 맥락·정적 지식 & $\bullet$ & $\circ$ &  &  &  & $\bullet$ &  & $\circ$ & &  \\
& \ch{C10} 시간 유효성 &  &  &  &  &  &  & $\bullet$ & $\circ$ & $\bullet$&  \\
& \ch{C11} 라벨 부족·미국 특수성 & $\bullet$ &  & $\bullet$ &  & $\bullet$ &  &  &  & &  \\
& \ch{C12} 오류 전파·설명 가능성 &  &  &  &  &  & $\circ$ & $\bullet$ & $\bullet$ & $\bullet$& $\circ$ \\
& \ch{C13} LLM 형식 오류·환각 &  &  &  &  & $\bullet$ &  & $\bullet$ &  & &  \\
& \ch{C14} 지역 전이 (억양·용어) & $\bullet$ &  & $\circ$ &  & $\circ$ & $\circ$ &  &  & & $\bullet$ \\
& \ch{C15} 데이터 이용 제약 &  & $\circ$ &  &  & $\bullet$ &  &  &  &  & $\bullet$ \\
\bottomrule
\end{tabular}
\end{center}

%==========================================================
\section{샘플 실행 예 (입출력 인스턴스)}\label{sec:example}
%==========================================================
샘플의 첫 교환(46.9--59.0\,s)을 각 모듈의 입출력으로 전개한다. 여기서는 VAD 미사용 시 \texttt{Alpha 2398}로 오인식된 경우를 가정한다.

\begin{center}\small
\begin{tabularx}{\textwidth}{@{}l>{\raggedright\arraybackslash}X@{}}
\toprule
모듈 & 출력 인스턴스\\\midrule
M1a & $s_1=[46.9, 53.4)$, $s_2=[54.0, 59.0)$. 이 앞의 0--46.9\,s 무음은 제거된다. VAD 없이 돌렸을 때는 이 구간에서 \texttt{"Thanks for watching."}이 환각으로 생성되었다 (\ch{C1}).\\
M1b & $w_1=$\,\texttt{"Alpha 2398, turn right at Juliet, cross runway two two right, and continue via Alpha to your left"}\\
& $w_2=$\,\texttt{"All right, right turn Juliet, cross runway two two right at Juliet, Alpha to your left, Delta 2398"}\\
M2a & $r_1=\ATCO$ (콜사인이 문장 앞), $r_2=\PILOT$ (콜사인이 문장 뒤)\\
M2b & $u_1=(\hat c_1, I_1)$, $\hat c_1=\texttt{"Alpha 2398"}$,\\
& $I_1=\{(\cmd{TURN},-,-,\cmd{RIGHT}{\wedge}\cmd{AT}{:}\texttt{Juliet},\texttt{Juliet}),$\\
& $\phantom{I_1=\{}(\cmd{CROSS\_RUNWAY},-,-,-,\texttt{"two two right"}),\ (\cmd{TAXI\_VIA},-,-,-,\texttt{Alpha})\}$\\
& $u_2=(\texttt{"Delta 2398"},\{\text{같은 3개 개념}\})$\\
M3 & $\nu(\texttt{"two two right"})=\cmd{22R}$, $\ell(\cdot)=\cmd{Runway\_22R}$\\
& $\hat c_1$의 \texttt{Alpha}는 telephony에 없고 숫자가 일치하며 $\mathrm{ctx}$가 높음 $\Rightarrow f_1=f_2=\cmd{DAL2398}$, $\mathcal{A}_1\ni\cmd{ASR\_CORRECTED}$ (\ch{C4})\\
M4 & $\cmd{Transmission}_2\xrightarrow{\cmd{READBACK\_OF}}\cmd{Transmission}_1$ ($t^s_2-t^e_1=0.6\,\mathrm{s}\le\Delta_{\mathrm{rb}}$). $\tau(\cmd{CROSS\_RUNWAY}) = [46.9, t^{\mathrm{to}})$\\
M5 & $\mathrm{Req}=\{\cmd{CROSS\_RUNWAY}\}$, $\mathrm{Miss}=\mathrm{Conf}=\emptyset\Rightarrow\cmd{MATCH}$. 단어 단위로는 순서·표현이 다르지만 개념 단위로는 일치한다 (\ch{C6}). 근거 $\pi$에는 \cmd{ASR\_CORRECTED} 표시가 포함된다.\\
M6 & $q_2(\cmd{DAL2398}) = \langle\cmd{TURN}\prec\cmd{CROSS\_RUNWAY}\prec\cmd{TAXI\_VIA}\rangle$\\
\bottomrule
\end{tabularx}
\end{center}

%==========================================================
\section{RQ4 일반화 검증 설계 (한국·일본 교신)}\label{sec:rq4}
%==========================================================
\textbf{팀 합의}: 시스템은 미국 교신 기준으로만 구축하고, 한국·일본 교신에는 \emph{추가 학습 없이} 적용해 검증만 한다.

\textbf{[해결 문제]} \ch{C14} 지역 전이.

\textbf{[왜 학습이 필요 없는가]} 시스템은 사전학습 STT(\md{M1b}), LLM few-shot(\md{M2b}), 규칙(\md{M5})으로 구성되어 학습 단계가 없다. 지역에 따라 달라지는 것은 \md{M0}의 자원뿐이다. 따라서 시스템을 지역 $\rho$에 대해 다음처럼 쓸 수 있다.
\begin{equation}
F_{\rho} = F\big(\,\cdot\,;\ \Ont_{\rho},\ \KS^{\rho},\ \Vdom^{\rho},\ \mathcal{E}_{\rho}\big),\qquad \rho\in\{\mathrm{US},\mathrm{KR},\mathrm{JP}\}.
\end{equation}
모듈 구현, 규칙의 형태 $\mathcal{R}$, LLM 백엔드는 모든 $\rho$에서 동일하다. 단, 복창 필수 집합 $\Treq^{\rho}$는 지역 규정(FAA 대 ICAO Doc 4444)에 따라 달라지는 \textbf{지역 자원}으로 취급한다.

\textbf{[실험 조건]} 대상 지역의 교신 집합을 $D_\rho$라 할 때 두 조건을 비교한다.
\begin{align}
\text{(a) 그대로 적용:}&\quad \hat{\G}^{(a)}_\rho = F_{\mathrm{US}}(D_\rho),\\
\text{(b) 지역 자원만 교체:}&\quad \hat{\G}^{(b)}_\rho = F\big(D_\rho;\ \Ont_{\mathrm{US}}[\Treq^{\rho}],\ \KS^{\rho},\ \Vdom^{\rho},\ \mathcal{E}_{\mathrm{US}}\big).
\end{align}
(b)에서 바꾸는 것은 공항 1곳(KR: \cmd{RKSI}, JP: \cmd{RJTT})의 정적 지식, 프롬프트 어휘, 그리고 $\KR$(\texttt{rules.json})의 복창 규칙(ICAO)뿐이고 온톨로지(\texttt{ontology.json})는 그대로다. ICAO·국가 규정 원문은 아직 입력 자료에 없어 RQ4 단계에서 확보·검증한다. 데이터는 OurAirports와 FAA JO 7340.2 콜사인 표 \cite{faa7340}를 그대로 재사용한다.

\textbf{[측정]} 국가별 송신 $n_\rho\approx 15$--$20$개에 대해 사람이 정답 $(f^*_i, I^*_i)$를 단다.
\begin{align}
\mathrm{CSR}_\rho &= \frac{1}{n_\rho}\sum_{i}\mathbb{1}\big[f_i = f^*_i\big] &&\text{(콜사인 인식률)}\\
\mathrm{CmdF1}_\rho &= \mathrm{F1}\Big(\textstyle\bigcup_i \tilde I_i,\ \bigcup_i I^*_i\Big) &&\text{(개념 단위 명령 인식, $m(\cdot,\cdot)$ 기준 일치)}\\
m^\dagger(i) &= \text{첫 } k\in(\mathrm{M1},\mathrm{M2a},\mathrm{M2b},\mathrm{M3}) \text{ s.t. } \mathrm{out}_{k}(i)\neq \mathrm{gold}_{k}(i) &&\text{(처음 실패한 단계)}
\end{align}
계산은 \md{M7}이 담당한다. $m^\dagger$의 분포로 \textbf{모듈별 실패 원인 표}를 만든다. $\mathrm{lang}_i\neq\texttt{en}$인 송신은 $m^\dagger=\mathrm{M1}$의 별도 범주(\cmd{UNSUPPORTED\_LANG})로 센다. 예: 억양 때문에 M1(STT)에서 실패했는지, 용어 차이 때문에 M2b에서 실패했는지, 공항 지식이 없어 M3에서 실패했는지. 이 측정은 KG 자체의 평가가 아니라 \emph{파이프라인 진단}이므로, KG 평가를 생략한다는 범위 결정과 충돌하지 않는다.

\textbf{[해석]}
\begin{itemize}
\item (b)에서 (a) 대비 회복되면: ``지역 자원만 교체하면 확장된다''는 모듈화 설계(\md{M0} 분리)가 실험으로 뒷받침된다.
\item 회복되지 않는 실패($m^\dagger=\mathrm{M1}$의 억양 오류, 자국어 혼용)는 지역 데이터 미세조정이 필요한 부분으로, 향후 과제의 근거가 된다 \cite{usatc2025,chen2023}.
\end{itemize}

\textbf{[타당성 위협]}
\begin{itemize}
\item \textbf{KR-1안(VATSIM)}: 지역 효과와 시뮬레이션 효과가 섞인다. 결과를 ``KR(모의)''로 표기하고, 가능하면 VATSIM 미국 교신 대조군으로 분리한다.
\item \textbf{표본 크기}: 국가별 15--20개라 통계적 일반화가 아니라 \emph{진단적 사례 분석}으로 보고한다.
\item \textbf{공항 차이}: KATL(지상) 대 RKSI/RJTT의 운영 차이가 섞일 수 있으므로, 가능한 한 같은 관제 위치(Ground)끼리 비교한다.
\end{itemize}

%==========================================================
\section{모듈 입출력 요약}
%==========================================================
\begin{center}\small
\begin{tabularx}{\textwidth}{@{}l>{\raggedright\arraybackslash}p{3.3cm}>{\raggedright\arraybackslash}p{3.6cm}>{\raggedright\arraybackslash}X@{}}
\toprule
모듈 & 입력 & 출력 & 구현 후보\\\midrule
M0 & $D_{\mathrm{apt}},D_{\mathrm{tel}},D_{\mathrm{fac}},D_{\mathrm{onto}},D_{\mathrm{reg}}$ & $\Ont,\ \KS,\ \KR,\ \Vdom$ (+ 설정 $\Cfg$) & Python, pdftotext (\texttt{kb/})\\
M1a & $\{(a^{\varphi},\Ctx^{\varphi})\}_{\varphi}$ & $\{(s_i,x_i,\Ctx_i)\}$ (UTC 병합) & ffmpeg, Silero VAD, 파일명 파서\\
M1b & $x_i,\Vdom,\mathcal{F}_{\mathrm{act}}$ & $(w_i,c_i,\mathrm{lang}_i)$ & faster-whisper large-v3 (RTX 4090에서 동작 확인)\\
M2a & $\{(w_i,\sigma_i)\}_{i\in\text{파일}}$ & $\{r_i\}$, $\rho:\sigma\mapsto r$ & LLM(화자 단위) + 규칙 다수결\\
M2b & $w_i,\Ont,\mathcal{E}$ & $u_i=(\hat c_i,I_i)$ & LLM(API 또는 로컬) + JSON Schema 구조화 출력\\
M3 & $u_i,\KS,\G_{<t_i},\Ctx_i$ & $\tilde u_i=(f_i,\tilde I_i,\mathcal{A}_i)$ & rapidfuzz, Double Metaphone, 규칙\\
M4 & $\G_{t_{i-1}},\tilde u_i,\Sigma,\KR,\Cfg$ & $\G_{t_i}$ & Neo4j + Cypher MERGE (Graphiti 참고)\\
M5 & $\G_t,\mathcal{R},\KR$ & $\Yset$ & Cypher 규칙 질의\\
M6 & $\G,t,q$ & $Q(\G,t,q)$ & Neo4j Browser / Streamlit\\
M7 & $\mathcal{J}$, 소량 정답 & $\mathrm{CSR}_\rho,\mathrm{CmdF1}_\rho,P_\rho(m^\dagger)$ & Python 로그 + 간단한 주석 시트\\
\bottomrule
\end{tabularx}
\end{center}

\paragraph{범위 밖 (명시).} TKG 임베딩·완성·예측, 실시간 스트리밍, 한국·일본 데이터로 학습·미세조정, 자국어 발화 인식(\cmd{UNSUPPORTED\_LANG}으로 기록만 함), ADS-B 실시간 연동은 제외한다. ADS-B는 \md{M1b}/\md{M3}의 $\mathcal{F}_{\mathrm{act}}$를 대체하는 확장점으로만 남긴다 \cite{nigmatulina2022}. 정량적 KG 평가는 교수 지시에 따라 데모로 대체한다.

%==========================================================
\begin{thebibliography}{99}\small
\bibitem{faa7110} FAA. \emph{Order JO 7110.65BB, Air Traffic Control}, effective 2025-02-20. §2-4-3 Pilot Acknowledgment/Read Back, §3-7-2 Taxi and Ground Movement Operations. \url{https://www.faa.gov/documentLibrary/media/Order/7110.65BB_Basic_dtd_2-20-25.pdf}
\bibitem{dlr2025} Helmke, H., Kleinert, M., Ohneiser, O. \emph{Ontology for Air Traffic Control Command Annotation -- Description}, v3.00. DLR-IB-FL-BS-2025-107, 2025-07-04, with OntologyATCCommandsTable\_v20250704. \url{https://elib.dlr.de/214852}
\bibitem{faa7340} FAA. \emph{Order JO 7340.2, Contractions} (항공사 telephony 및 3자리 코드).
\bibitem{eurocontrol} Van Es, G. (NLR for EUROCONTROL). \emph{Air-Ground Communication Safety Study: Causes and Recommendations}, Ed.~1.1, 2005.
\bibitem{cardosi94} Cardosi, K.~M. \emph{An Analysis of Tower (Local) Controller--Pilot Voice Communications}. DOT/FAA/RD-94/15, Volpe Center, 1994.
\bibitem{helmke2025loss} Helmke, H., Clarke, S.~S., Ohneiser, O., Kleinert, M., Shetty, S., Kalyanam, K. Automatic Speech Recognition and Understanding in Air Traffic Control -- How Much Information Is Lost Over the Atlantic Ocean? \emph{DASC}, 2025. NASA NTRS 20250001345.
\bibitem{helmke2018} Helmke, H. et al. Ontology for Transcription of ATC Speech Commands of SESAR 2020 Solution PJ.16-04. \emph{DASC}, 2018. doi:10.1109/DASC.2018.8569238
\bibitem{kleinert2021} Kleinert, M., Helmke, H., Shetty, S., Ohneiser, O., Ehr, H., Prasad, A., Motlicek, P., Harfmann, J. Automated Interpretation of Air Traffic Control Communication: The Journey from Spoken Words to a Deeper Understanding of the Meaning. \emph{DASC}, 2021. doi:10.1109/DASC52595.2021.9594387
\bibitem{helmke2021} Helmke, H. et al. Readback Error Detection by Automatic Speech Recognition to Increase ATM Safety. \emph{14th USA/Europe ATM R\&D Seminar}, 2021.
\bibitem{chen2023} Chen, S., Helmke, H., Tarakan, R.~M., Ohneiser, O., Kopald, H., Kleinert, M. Effects of Language Ontology on Transatlantic Automatic Speech Understanding Research Collaboration in the ATM Domain. \emph{Aerospace} 10(6):526, 2023.
\bibitem{keller2017} Keller, R.~M. \emph{The NASA Air Traffic Management Ontology}. NASA/TM-2017-219526, 2017.
\bibitem{keller2019} Keller, R.~M. Building a Knowledge Graph for the Air Traffic Management Community. \emph{Companion Proc. WWW}, 2019. doi:10.1145/3308560.3317706
\bibitem{atco2} Zuluaga-Gomez, J. et al. ATCO2 corpus: A Large-Scale Dataset for Research on ASR and NLU of ATC Communications. \emph{DMLR}, 2024. arXiv:2211.04054
\bibitem{bertraffic} Zuluaga-Gomez, J. et al. BERTraffic: BERT-based Joint Speaker Role and Speaker Change Detection for ATC Communications. \emph{IEEE SLT}, 2023.
\bibitem{nigmatulina2022} Nigmatulina, I., Zuluaga-Gomez, J., Prasad, A., Sarfjoo, S.~S., Motlicek, P. A Two-Step Approach to Leverage Contextual Data: Speech Recognition in Air-Traffic Communications. \emph{ICASSP}, 2022.
\bibitem{whisper} Radford, A. et al. Robust Speech Recognition via Large-Scale Weak Supervision. \emph{ICML}, 2023.
\bibitem{vandoorn} van Doorn, J. et al. Whisper-ATC (Whisper fine-tuned on ATCO2/ATCOSIM). TU Delft. (정확한 제목·학회 확인 필요)
\bibitem{usatc2025} \emph{Fine-Tuning Whisper for American English Air Traffic Control Speech Recognition: A Data-Efficient Pipeline}. Research Square preprint. (저자 확인 필요)
\bibitem{koenecke2024} Koenecke, A., Choi, A.~S.~G., Mei, K.~X., Schellmann, H., Sloane, M. Careless Whisper: Speech-to-Text Hallucination Harms. \emph{ACM FAccT}, 2024. doi:10.1145/3630106.3658996
\bibitem{mekerishvili2025} Mekerishvili, A.~M., Sun, J., Jonk, P., de Vries, V. Structured Command Extraction from ATC Communications Using Open and Fine-Tuned Language Models. \emph{SESAR Innovation Days}, 2025.
\bibitem{llmkgc2025} LLM-empowered Knowledge Graph Construction: A Survey. arXiv:2510.20345, 2025.
\bibitem{feng2024} Feng, X., Wu, X., Meng, H. Ontology-grounded Automatic Knowledge Graph Construction by LLM under Wikidata Schema. \emph{HI-AI@KDD}, 2024. arXiv:2412.20942
\bibitem{edc2024} Zhang, B., Soh, H. Extract, Define, Canonicalize: An LLM-based Framework for Knowledge Graph Construction. \emph{EMNLP}, 2024.
\bibitem{itext2kg} Lairgi, Y., Moncla, L., Cazabet, R., Benabdeslem, K., Cl\'eau, P. iText2KG: Incremental Knowledge Graphs Construction Using Large Language Models. \emph{WISE}, 2024. arXiv:2409.03284
\bibitem{atom2025} Lairgi, Y. et al. ATOM: AdapTive and OptiMized Dynamic Temporal Knowledge Graph Construction Using LLMs. arXiv:2510.22590, 2025.
\bibitem{zep2025} Rasmussen, P., Paliychuk, P., Beauvais, T., Ryan, J., Chalef, D. Zep: A Temporal Knowledge Graph Architecture for Agent Memory. arXiv:2501.13956, 2025.
\bibitem{graphrag} Edge, D. et al. From Local to Global: A Graph RAG Approach to Query-Focused Summarization. arXiv:2404.16130, 2024.
\bibitem{willard2023} Willard, B.~T., Louf, R. Efficient Guided Generation for Large Language Models. arXiv:2307.09702, 2023.
\bibitem{naryw3c} Noy, N., Rector, A. (eds.) \emph{Defining N-ary Relations on the Semantic Web}. W3C Working Group Note, 2006.
\bibitem{shacl} Knublauch, H., Kontokostas, D. (eds.) \emph{Shapes Constraint Language (SHACL)}. W3C Recommendation, 2017.
\bibitem{owltime} Cox, S., Little, C. (eds.) \emph{Time Ontology in OWL}. W3C Recommendation, 2017.
\bibitem{allen1983} Allen, J.~F. Maintaining Knowledge about Temporal Intervals. \emph{Commun. ACM} 26(11):832--843, 1983.
\bibitem{knez2023} Knez, T., \v{Z}itnik, S. Event-Centric Temporal Knowledge Graph Construction: A Survey. \emph{Mathematics} 11(23):4852, 2023.
\bibitem{cai2024tkg} Cai, L. et al. A Survey on Temporal Knowledge Graph: Representation Learning and Applications. arXiv:2403.04782, 2024.
\bibitem{pan2024} Pan, S., Luo, L., Wang, Y., Chen, C., Wang, J., Wu, X. Unifying Large Language Models and Knowledge Graphs: A Roadmap. \emph{IEEE TKDE} 36(7):3580--3599, 2024.
\bibitem{endsley1995} Endsley, M.~R. Toward a Theory of Situation Awareness in Dynamic Systems. \emph{Human Factors} 37(1):32--64, 1995.
\bibitem{hfliveatc} nyuuzyou. \emph{LiveATC Recordings (Partial 2024-08-26)}. Hugging Face Datasets, \url{https://huggingface.co/datasets/nyuuzyou/liveatc}. LiveATC.net Terms of Use 적용(비상업, 재배포 허가 필요).
\bibitem{icao4444} ICAO. \emph{Doc 4444, Procedures for Air Navigation Services -- Air Traffic Management (PANS-ATM)}. (복창 요건 조항 번호는 최신판에서 확인 필요)
\bibitem{atc0} Godfrey, J.~J. \emph{Air Traffic Control Complete} (ATC0). LDC94S14A, Linguistic Data Consortium, 1994. doi:10.35111/2bg6-nn53
\bibitem{tartan2025} Patrikar, J. et al. Image, Speech, and ADS-B Trajectory Datasets for Terminal Airspace Operations (TartanAviation). \emph{Scientific Data}, 2025. doi:10.1038/s41597-025-04775-6
\bibitem{uwbatcc} \v{S}m\'idl, L., \v{S}vec, J., Tihelka, D., Matou\v{s}ek, J., Romportl, J., Ircing, P. Air Traffic Control Communication (ATCC) Speech Corpora and Their Use for ASR and TTS Development. \emph{Language Resources and Evaluation} 53, 2019.
\bibitem{atcosim} Hofbauer, K., Petrik, S., Hering, H. The ATCOSIM Corpus of Non-Prompted Clean Air Traffic Control Speech. \emph{LREC}, 2008.
\bibitem{airbusatc} Pellegrini, T. et al. The Airbus Air Traffic Control Speech Recognition 2018 Challenge: Towards ATC Automatic Transcription and Call Sign Detection. \emph{Interspeech}, 2019. arXiv:1810.12614
\bibitem{atcspeech} Yang, B. et al. ATCSpeech: A Multilingual Pilot-Controller Speech Corpus from Real Air Traffic Control Environment. \emph{Interspeech}, 2020. arXiv:1911.11365
\bibitem{youtubeatc} W\"ustenbecker, N., Ohneiser, O., Kleinert, M. Can YouTube Stream Recordings Improve Speech Recognition for Air Traffic Control? \emph{Journal of Open Aviation Science} 4, 2026.
\bibitem{zuluaga2021} Zuluaga-Gomez, J., Nigmatulina, I., Prasad, A., Motlicek, P., Vesel\'y, K., Kocour, M., Sz\H{o}ke, I. Contextual Semi-Supervised Learning: An Approach to Leverage Air-Surveillance and Untranscribed ATC Data in ASR Systems. \emph{Interspeech}, 2021.
\bibitem{scope2026} Deng, Q. et al. SCOPE: A Lightweight-training LLM Framework for Air Traffic Control Readback Monitoring. arXiv:2605.29543, 2026.
\end{thebibliography}

\end{document}
